\chapter{動態規劃}
\label{ch:dp}

% 本章的圖共用的樣式（名稱都加 dp- 前綴，避免和別章撞名）
\tikzset{
    dp-call/.style={rounded corners=2pt, draw=black!60, fill=white, minimum width=0.85cm, minimum height=0.5cm, inner sep=1.5pt, font=\small},
    dp-look/.style={dp-call, draw=memStroke, thick, fill=memFill},
    dp-gone/.style={dp-call, draw=black!25, text=black!35, fill=black!4, dashed},
    dp-base/.style={dp-call, fill=black!6},
    dp-leaf/.style={dp-call, draw=coreStroke, fill=coreFill},
    dp-e/.style={draw=black!55},
    dp-gonee/.style={draw=black!25, dashed},
    dp-arr/.style={-{Stealth[length=1.7mm]}, draw=black!55},
    dp-hd/.style={font=\small, text=black!70},
    dp-c/.style={draw=black!40, minimum width=0.95cm, minimum height=0.62cm, inner sep=0pt, font=\small},
    dp-c0/.style={dp-c, fill=black!5, text=black!55},
    dp-cd0/.style={dp-c},
    dp-cd1/.style={dp-c, fill=coreFill},
    dp-cpath/.style={dp-c, draw=memStroke, thick, fill=memFill},
    dp-chit/.style={dp-c, draw=optStroke, thick, fill=optFill},
    dp-key/.style={rounded corners=2pt, draw=black!65, fill=white, minimum height=0.5cm, inner sep=2pt, font=\small},
    dp-ext/.style={rectangle, draw=black!45, fill=black!6, minimum size=0.42cm, inner sep=1pt, font=\scriptsize},
}

系上的程式作業允許同學在截止前重交。 助教手上有同一位同學交的第一版和第二版，想知
道第二版到底改了多少：兩個版本裡有多少行可以原封不動地對上，對不上的才算改過。 從
頭一行一行往下對很容易對錯位，第二版只要在前面插了一行，後面每一行就都對不上了；可
是把「哪些行要對上」的選法全部試一遍，三百行的程式要試多少種？ 往下讀之前，先想一
想：如果已經知道兩個版本「前面一段」最多能對上幾行，能不能拿這個答案去算更長的一
段？

這一章講\textbf{動態規劃}（dynamic programming，常簡稱 DP）。 它和第~17~章的分治一樣，
把問題拆成較小的子問題；不同的是，分治拆出來的子問題互不相干，動態規劃處理的是子問題
\textbf{會重複出現}的情形。 做法是每個子問題只解一次，把答案記在一張表裡，之後要用時直
接查表。 名稱裡的 programming 是二十世紀中期「規劃、排表」的意思，和寫程式無關。

這一章先用一個計數的例子看出重複的子問題，比較由上而下的記憶化和由下而上的填表，再整
理一個最佳化問題要能用動態規劃的條件。 接著五節是經典問題：矩陣鏈乘、最長共同子序列、
最佳二元搜尋樹、0/1 背包，以及編輯距離與最長遞增子序列，每一個都照考卷的寫法把表格填完、
再把最佳解還原出來。 最後用動態規劃的眼光看所有點對最短路徑，完整的內容在第~25~章。

這一章用到第~\ref{sec:ana-input-size}~節的擬多項式時間、第~2~章的遞迴樹、第~7~章的卡特蘭數
與延伸二元樹、第~9~章的二元搜尋樹，以及第~\ref{sec:dc-bsearch}~節的 \proc{Lower-Bound}。 C 程式用到《C
程式設計基礎》第~9~章把二維表攤平成一維陣列的寫法、第~15~章的記憶化和第~19~章的
\texttt{malloc}、\texttt{free}，用到時會先在框裡複習。

\section{動態規劃的要素}
\label{sec:dp-elements}

\subsection{重疊的子問題：數呼叫次數}
\label{sec:dp-overlap}

從 $n$ 個人裡選 $k$ 個人，方法數是二項式係數 $\binom{n}{k}$。 把第 $n$ 個人分開看：選
他的話，剩下要從 $n - 1$ 人裡選 $k - 1$ 人；不選他的話，要從 $n - 1$ 人裡選 $k$ 人。
這就是巴斯卡公式
\[
\binom{n}{k} = \binom{n-1}{k-1} + \binom{n-1}{k}, \qquad
\binom{n}{0} = \binom{n}{n} = 1 .
\]
照公式直接寫成遞迴：

\begin{pseudo}[label=alg:dp-plain-binom]{照巴斯卡公式直接遞迴}
\Function{Plain-Binom}{$n, k$}
    \If{$k = 0$ \OR $k = n$}
        \State \Return $1$
    \EndIf
    \State \Return $\Call{Plain-Binom}{n - 1, k - 1} + \Call{Plain-Binom}{n - 1, k}$
\EndFunction
\end{pseudo}

下圖是 $\proc{Plain-Binom}(5, 2)$ 的遞迴樹，每個節點寫的是那一次呼叫的參數 $n, k$。
$\binom{5}{2} = 10$，樹上剛好有 $10$ 個回傳 $1$ 的葉節點，因為答案就是這些 $1$ 加起來
的。 同一組參數出現了好幾次：$(3, 1)$ 出現 $2$ 次，$(2, 1)$ 出現 $3$ 次，每一次都從頭
重算一遍。

\begin{center}
\begin{minipage}{\linewidth}
\centering
\begin{tikzpicture}[x=1.32cm, y=1.05cm]
\node[dp-call] (b0) at (3.69,0) {$5,2$};
\node[dp-call] (b1) at (0.88,-1) {$4,1$};
\node[dp-call] (b2) at (0.00,-2) {$3,0$};
\node[dp-call] (b3) at (1.75,-2) {$3,1$};
\node[dp-call] (b4) at (1.00,-3) {$2,0$};
\node[dp-call] (b5) at (2.50,-3) {$2,1$};
\node[dp-call] (b6) at (2.00,-4) {$1,0$};
\node[dp-call] (b7) at (3.00,-4) {$1,1$};
\node[dp-call] (b8) at (6.50,-1) {$4,2$};
\node[dp-look] (b9) at (4.75,-2) {$3,1$};
\node[dp-gone] (b10) at (4.00,-3) {$2,0$};
\node[dp-gone] (b11) at (5.50,-3) {$2,1$};
\node[dp-gone] (b12) at (5.00,-4) {$1,0$};
\node[dp-gone] (b13) at (6.00,-4) {$1,1$};
\node[dp-call] (b14) at (8.25,-2) {$3,2$};
\node[dp-look] (b15) at (7.50,-3) {$2,1$};
\node[dp-gone] (b16) at (7.00,-4) {$1,0$};
\node[dp-gone] (b17) at (8.00,-4) {$1,1$};
\node[dp-call] (b18) at (9.00,-3) {$2,2$};
\draw[dp-e] (b0) -- (b1);
\draw[dp-e] (b1) -- (b2);
\draw[dp-e] (b1) -- (b3);
\draw[dp-e] (b3) -- (b4);
\draw[dp-e] (b3) -- (b5);
\draw[dp-e] (b5) -- (b6);
\draw[dp-e] (b5) -- (b7);
\draw[dp-e] (b0) -- (b8);
\draw[dp-e] (b8) -- (b9);
\draw[dp-gonee] (b9) -- (b10);
\draw[dp-gonee] (b9) -- (b11);
\draw[dp-gonee] (b11) -- (b12);
\draw[dp-gonee] (b11) -- (b13);
\draw[dp-e] (b8) -- (b14);
\draw[dp-e] (b14) -- (b15);
\draw[dp-gonee] (b15) -- (b16);
\draw[dp-gonee] (b15) -- (b17);
\draw[dp-e] (b14) -- (b18);
\end{tikzpicture}
\diagcaption{$\proc{Plain-Binom}(5, 2)$ 的遞迴樹，共 $19$ 次呼叫。 依「先左後右」的順序執
行時，橘色的兩個節點是「這組參數之前算過了」的呼叫；虛線灰色的 $6$ 個節點是它們底下的
重算。 下一小節的記憶化會在橘色節點直接查表，灰色的呼叫都不會發生。}
\label{fig:dp-binom-tree}
\end{minipage}
\end{center}

\begin{theorem}
\label{thm:dp-plain-calls}
$0 \le k \le n$ 時，$\proc{Plain-Binom}(n, k)$ 一共被呼叫 $2\binom{n}{k} - 1$ 次。
\end{theorem}

\begin{proof}
令呼叫次數為 $T(n, k)$，對 $n$ 用歸納法。 $k = 0$ 或 $k = n$ 時只有一次呼叫，
$T = 1 = 2 \cdot 1 - 1$。 其他情況，這一次呼叫加上兩個子呼叫：
\[
T(n, k) = 1 + T(n-1, k-1) + T(n-1, k)
= 1 + \left(2\tbinom{n-1}{k-1} - 1\right) + \left(2\tbinom{n-1}{k} - 1\right)
= 2\tbinom{n}{k} - 1 .
\]
最後一步用了巴斯卡公式。
\end{proof}

這個次數是指數級的：$\binom{2m}{m}$ 是 $\binom{2m}{0}, \binom{2m}{1}, \ldots, \binom{2m}{2m}$
這 $2m + 1$ 個數裡最大的，而它們加起來是 $2^{2m} = 4^m$，所以 $\binom{2m}{m} \ge 4^m / (2m + 1)$。 可是會出現的參數組合 $(i, j)$ 只有大約
$n \cdot k$ 種。 呼叫次數和「不同的子問題個數」差這麼多，表示同一個子問題被算了很多次。
這種情形叫\textbf{重疊的子問題}（overlapping subproblems）。 第~14~章的合併排序不一樣：
每次遞迴處理的都是陣列裡新的一段，沒有一段會被排兩次，所以記表格也省不了什麼。

第~2~章的鋪地磚程式 \texttt{tile(n)} 也是同一種情形，那裡算過呼叫次數是 $\Theta(2^n)$，
而不同的參數只有 $n + 1$ 個。

\subsection{由上而下：記憶化}
\label{sec:dp-memo}

最直接的修法是在遞迴外面加一張表：第一次算出 $\binom{i}{j}$ 時寫進表裡，之後再被問到
就直接回傳。 這叫\textbf{記憶化}（memoization），《C 程式設計基礎》第~15~章 15.9 節用費氏
數列示範過。 表 $\mathit{memo}[0..n, 0..k]$ 一開始每一格都是 \NIL。

\begin{pseudo}[label=alg:dp-memo-binom]{記憶化的二項式係數}
\Function{Memo-Binom}{$n, k$}
    \If{$k = 0$ \OR $k = n$}
        \State \Return $1$
    \EndIf
    \If{$\mathit{memo}[n, k] = \NIL$}
        \State $\mathit{memo}[n, k] \gets \Call{Memo-Binom}{n - 1, k - 1} + \Call{Memo-Binom}{n - 1, k}$
    \EndIf
    \State \Return $\mathit{memo}[n, k]$
\EndFunction
\end{pseudo}

要算幾次？ 從 $(n, k)$ 往下走，每一步不是 $(i, j) \to (i-1, j-1)$（「選」）就是
$(i, j) \to (i-1, j)$（「不選」）。 「選」最多走 $k$ 次、「不選」最多走 $n - k$ 次，而
一個子問題由「已經走了幾次選、幾次不選」決定。 會走到的子問題裡，不是邊界的（$j \ge 1$
而且 $j < i$）正是「選」走了不到 $k$ 次、「不選」走了不到 $n - k$ 次的那些，共
$k(n - k)$ 個。 每個這樣的子問題在第一次被問到時展開一次，發出 $2$ 個呼叫，其他呼叫都在
第 3 行或第 6 行直接回傳，所以
\[
\proc{Memo-Binom}(n, k) \text{ 的呼叫次數} = 1 + 2k(n - k) .
\]
$n = 5$、$k = 2$ 時是 $13$，正好是圖~\ref{fig:dp-binom-tree} 裡白色和橘色節點的個數。
時間從指數變成 $O(nk)$。

把每個不同的子問題畫成一個節點，從每個子問題畫箭頭指向它直接用到的子問題，得到
\textbf{子問題圖}（subproblem graph）。 遞迴樹把同一個子問題畫很多次，子問題圖只畫一
次。 記憶化的呼叫次數等於 $1$ 加上子問題圖的邊數。

\begin{center}
\begin{minipage}{\linewidth}
\centering
\begin{tikzpicture}[x=1.45cm, y=1.0cm]
\node[dp-base] (g10) at (2.00,-4) {$1,0$};
\node[dp-base] (g11) at (3.00,-4) {$1,1$};
\node[dp-base] (g20) at (1.50,-3) {$2,0$};
\node[dp-call] (g21) at (2.50,-3) {$2,1$};
\node[dp-base] (g22) at (3.50,-3) {$2,2$};
\node[dp-base] (g30) at (1.00,-2) {$3,0$};
\node[dp-call] (g31) at (2.00,-2) {$3,1$};
\node[dp-call] (g32) at (3.00,-2) {$3,2$};
\node[dp-call] (g41) at (1.50,-1) {$4,1$};
\node[dp-call] (g42) at (2.50,-1) {$4,2$};
\node[dp-call] (g52) at (2.00,0) {$5,2$};
\draw[dp-arr] (g21) -- (g10);
\draw[dp-arr] (g21) -- (g11);
\draw[dp-arr] (g31) -- (g20);
\draw[dp-arr] (g31) -- (g21);
\draw[dp-arr] (g32) -- (g21);
\draw[dp-arr] (g32) -- (g22);
\draw[dp-arr] (g41) -- (g30);
\draw[dp-arr] (g41) -- (g31);
\draw[dp-arr] (g42) -- (g31);
\draw[dp-arr] (g42) -- (g32);
\draw[dp-arr] (g52) -- (g41);
\draw[dp-arr] (g52) -- (g42);
\end{tikzpicture}
\diagcaption{$\binom{5}{2}$ 的子問題圖：$11$ 個不同的子問題，灰底的是邊界（直接回傳
$1$），$6$ 個白色的各有兩條邊，共 $12$ 條邊。 它沒有環，是一個有向無環圖。}
\label{fig:dp-binom-dag}
\end{minipage}
\end{center}

\subsection{由下而上：填表}
\label{sec:dp-bottomup}

另一種做法完全不遞迴：先想清楚子問題之間誰依賴誰，再從最小的子問題開始，依序把表填滿。
$\binom{i}{j}$ 只用到第 $i - 1$ 列，所以一列一列由上往下填就好。

\begin{pseudo}[label=alg:dp-binom-table]{由下而上填巴斯卡三角形}
\Function{Binom-Table}{$n, k$}
    \For{$i \gets 0$ \To $n$}
        \For{$j \gets 0$ \To $\min(i, k)$}
            \If{$j = 0$ \OR $j = i$}
                \State $C[i, j] \gets 1$
            \Else
                \State $C[i, j] \gets C[i - 1, j - 1] + C[i - 1, j]$
            \EndIf
        \EndFor
    \EndFor
    \State \Return $C[n, k]$
\EndFunction
\end{pseudo}

第 7 行執行時，第 $i - 1$ 列已經填好了，這是由下而上的關鍵：\textbf{填某一格之前，它用到
的格子都要先填好}。 用子問題圖來說，就是箭頭指到的節點要排在前面，也就是把子問題圖做
\textbf{拓樸排序}（第~22~章）再倒過來。 這一章的問題，順序都可以直接看出來：依列號、依
區間長度、依前綴長度。

第 7 行只用到上一列，所以整張表可以換成一列 $\mathit{row}[0..k]$，每一列就地更新。 要注意
更新的方向：新的 $\mathit{row}[j]$ 要用\textbf{舊的} $\mathit{row}[j - 1]$，所以 $j$ 要由
大往小跑，否則 $\mathit{row}[j - 1]$ 已經先被改成新的值了。 第~\ref{sec:dp-knapsack}~節的
0/1 背包會再遇到同一件事。

看實際的數字：$\binom{24}{12} = 2{,}704{,}156$，\proc{Plain-Binom} 要呼叫
$2\binom{24}{12} - 1 = 5{,}408{,}311$ 次，\proc{Memo-Binom} 只要 $1 + 2 \cdot 12 \cdot 12 = 289$
次。 \proc{Binom-Table} 和一列滾動的版本，第 $i$ 列都做 $\min(i - 1, k)$ 次加法，共 $nk - k(k+1)/2 = 210$ 次；它也算了一
些用不到的格子，例如算 $\binom{24}{12}$ 根本用不到 $\binom{20}{1}$，但沒有遞迴，常數很小。

\subsection{最佳子結構}
\label{sec:dp-substructure}

上面是計數問題。 動態規劃更常用在\textbf{最佳化問題}（optimization problem）：每個可行解
有一個值，要找值最小（或最大）的解。 這時除了子問題要重疊，還需要第二個條件。

\begin{definition}[最佳子結構]
\label{def:dp-substructure}
一個問題有\textbf{最佳子結構}（optimal substructure），是說它的最佳解裡面包含的子問題
的解，也都是那些子問題的最佳解。
\end{definition}

有了這個性質，才能寫出「最佳值 $=$ 在幾種選擇中，取（子問題的最佳值 $+$ 這次選擇的成本）
最小的一個」這種遞迴式。 證明一個問題有最佳子結構，幾乎都用同一個手法，叫做
\textbf{剪下貼上}（cut-and-paste）：假設最佳解裡的某個子問題的解不是最佳的，就把它剪掉、
貼上那個子問題真正的最佳解，得到一個更好的整體解，和「原本是最佳解」矛盾。

\textbf{有最佳子結構的例子：最短路徑。} 在一張圖上，如果 $p$ 是從 $u$ 到 $v$ 邊數最少的
路徑，而 $w$ 是 $p$ 上的一個點，那麼 $p$ 從 $u$ 到 $w$ 的那一段 $p_1$ 也是 $u$ 到 $w$ 邊
數最少的路徑。 理由：如果有更短的 $p_1'$，把 $p_1$ 換成 $p_1'$，得到一條從 $u$ 到 $v$
的走法，比 $p$ 短。 這條走法可能經過同一個點兩次，但把中間繞回來的一圈去掉只會更短，
最後得到一條比 $p$ 短的路徑，矛盾。

\textbf{沒有最佳子結構的例子：最長簡單路徑。} 簡單路徑是不重複經過任何點的路徑。 看一個
三角形：三個點 $a, b, c$ 兩兩相連。 從 $a$ 到 $c$ 最長的簡單路徑是 $a \to b \to c$，有
$2$ 條邊。 它從 $a$ 到 $b$ 的那一段只有 $1$ 條邊，可是從 $a$ 到 $b$ 最長的簡單路徑是
$a \to c \to b$，有 $2$ 條邊。 最佳解的一段不是那一段的最佳解。 反過來把兩段的最佳解接起
來，$a \to c \to b$ 接上 $b \to a \to c$，會重複經過 $a$ 和 $c$，根本不是簡單路徑。

差別在於子問題是否\textbf{互相獨立}。 最短路徑的兩段各自用到的點不會互相限制（真的重複
了也能去掉）；最長簡單路徑的兩段卻會搶同一批點，一段用了某個點，另一段就不能用。 剪下
貼上的論證在這裡失效。 最長簡單路徑目前沒有已知的多項式時間演算法，第~28~章會講它和
NP 完備問題的關係。

\subsection{設計動態規劃的步驟}
\label{sec:dp-steps}

一個最佳化問題同時有最佳子結構和重疊的子問題時，可以照下面四步設計演算法：

\begin{enumerate}
\item 描述最佳解的結構：最佳解的「最後一步」或「第一步」有哪幾種選擇，做了某個選擇之後，
剩下的是哪個子問題。
\item 遞迴定義最佳值：寫出子問題的參數、遞迴式和邊界。
\item 計算最佳值：記憶化，或找出填表順序由下而上算。
\item 還原最佳解：填表時順便記下每一格是哪個選擇得到的，最後從答案那一格倒著走回去。
\end{enumerate}

只要最佳值的話，第 4 步可以省略。 執行時間通常是
\[
\text{子問題的個數} \times \text{每個子問題要比較的選擇數},
\]
再乘上每個選擇的成本（通常是常數）。 例如下一節的矩陣鏈乘有 $\Theta(n^2)$ 個子問題，每個
最多 $n - 1$ 種選擇，時間是 $O(n^3)$。

\begin{examnote}
考卷上的動態規劃題，答案通常要包含五件事，順序也照這樣寫：(1) 用一句完整的話定義表格的
一格，例如「$c[i, j]$ 是 $X$ 的前 $i$ 個字元和 $Y$ 的前 $j$ 個字元的最長共同子序列長
度」；(2) 遞迴式，連同邊界條件；(3) 填表的順序；(4) 答案在哪一格；(5) 時間和空間複雜度。
只寫出遞迴式、沒有寫那一格代表什麼，閱卷的人很難判斷你的遞迴式對不對。 填表題要把整張
表寫完，平手時照題目規定的規則選（本章各節都會寫明本書的規則）。
\end{examnote}

\subsection{由上而下還是由下而上}

兩種寫法算的是同一張表，漸近時間通常一樣。 記憶化照遞迴式直接寫，只算真的用得到的子問題，
但有函式呼叫的成本，遞迴也可能很深；由下而上要先想清楚填表順序，通常整張表都填，可是只有迴
圈、常數小，而且常常可以只留一兩列（滾動陣列）來省空間。
用得到的子問題只佔整張表的一小部分時，記憶化比較快，觀念題第~\ref{dp-q-sparse}~題有一個
例子。 遞迴深度可能到 $10^5$ 以上時（例如很長的字串），C 的呼叫堆疊可能不夠，要
寫成由下而上。 考卷上的填表題，幾乎都是由下而上的寫法。

\section{矩陣鏈乘}
\label{sec:dp-matrix}

\subsection{問題}

一個 $p \times q$ 的矩陣乘上一個 $q \times r$ 的矩陣，結果是 $p \times r$；用一般的三層迴圈
算，每一格要做 $q$ 次乘法，共 $pqr$ 次純量乘法。 這一節只數純量乘法的次數。

要算 $n$ 個矩陣的連乘 $A_1 A_2 \cdots A_n$，其中 $A_i$ 是 $p_{i-1} \times p_i$（相鄰兩個的
大小要接得上），所以整條鏈的大小由一個序列 $\langle p_0, p_1, \ldots, p_n \rangle$ 決定。
矩陣乘法有結合律，先乘哪兩個都得到同樣的結果，但次數可能差很多。 這一節的例子是 $6$ 個
矩陣，
\[
\langle p_0, \ldots, p_6 \rangle = \langle 6, 10, 3, 12, 5, 8, 4 \rangle ,
\]
也就是 $A_1$ 是 $6 \times 10$、$A_2$ 是 $10 \times 3$、…、$A_6$ 是 $8 \times 4$。 只看前三個：
\begin{itemize}
\item $(A_1 A_2) A_3$：先算 $6 \cdot 10 \cdot 3 = 180$ 次得到 $6 \times 3$，再乘 $A_3$ 要
$6 \cdot 3 \cdot 12 = 216$ 次，共 $396$ 次。
\item $A_1 (A_2 A_3)$：先算 $10 \cdot 3 \cdot 12 = 360$ 次得到 $10 \times 12$，再算
$6 \cdot 10 \cdot 12 = 720$ 次，共 $1080$ 次。
\end{itemize}
六個全部由左往右乘要 $1188$ 次，最差的加法要 $2472$ 次，等一下會算出最少是 $648$ 次。

\textbf{矩陣鏈乘問題}（matrix-chain multiplication problem）：給定 $\langle p_0, \ldots, p_n
\rangle$，找出一種加滿括號的方式，讓純量乘法的總次數最少。 注意輸出是\textbf{乘的順序}，
不是乘出來的矩陣。

括號的加法有幾種？ 第~7~章算過，$n + 1$ 個數連乘、加滿括號的方法有 $C_n$ 種（卡特蘭數），
所以 $n$ 個矩陣有 $C_{n-1}$ 種。 $C_{n-1} = \frac{1}{n}\binom{2n-2}{n-1}$ 是
$\Omega(4^n / n^{3/2})$，$n = 6$ 已經有 $42$ 種，$n = 20$ 有超過十億種，全部試一遍不可行。

\subsection{最佳子結構與遞迴式}

把 $A_i A_{i+1} \cdots A_j$ 簡寫成 $A_{i..j}$，它是 $p_{i-1} \times p_j$ 的矩陣。 $i < j$ 時，
任何一種加括號的方式，最後一次乘法一定是
\[
A_{i..j} = (A_{i..k}) \cdot (A_{k+1..j}), \qquad i \le k < j ,
\]
最後這一次的成本是 $p_{i-1} p_k p_j$。 如果最佳的乘法在 $k$ 分開，那麼它算 $A_{i..k}$ 的方
式一定也是 $A_{i..k}$ 最省的方式：假如有更省的，把那一段換掉，整體就更省了，這就是剪下
貼上。 右邊 $A_{k+1..j}$ 也一樣。 兩邊用到的矩陣不重疊，互不影響。

令 $m[i, j]$ 為算出 $A_{i..j}$ 最少需要的純量乘法次數。 我們不知道最佳的 $k$ 是哪一個，就每
一個都試：
\[
m[i, j] = \begin{cases}
0, & i = j, \\[2pt]
\displaystyle\min_{i \le k < j} \bigl\{\, m[i, k] + m[k+1, j] + p_{i-1} p_k p_j \,\bigr\}, & i < j .
\end{cases}
\]
另外用 $s[i, j]$ 記下取到最小值的那個 $k$，還原括號時要用。 整個問題的答案是 $m[1, n]$。

子問題是所有的 $1 \le i \le j \le n$，共 $\binom{n}{2} + n = \Theta(n^2)$ 個，而且大量重疊：
例如 $m[3, 5]$ 會被 $m[1, 5]$、$m[2, 5]$、$m[3, 6]$ 都用到。 直接照遞迴式遞迴會是指數時
間（觀念題第~\ref{dp-q-mc-rec}~題）。

\subsection{依鏈長填表}

$m[i, j]$ 用到的 $m[i, k]$ 和 $m[k+1, j]$ 都是比較短的鏈，所以依\textbf{鏈長}
$\ell = j - i + 1$ 由短到長填：先填 $\ell = 2$ 的所有格子，再填 $\ell = 3$，……，最後是
$\ell = n$ 的 $m[1, n]$。 畫成表格的話，$\ell$ 固定的格子排成一條和主對角線平行的斜線，所以
這種填法也叫「一條對角線一條對角線地填」。

\begin{pseudo}[label=alg:dp-mc-order]{矩陣鏈乘的填表}
\Function{Matrix-Chain-Order}{$p, n$}
    \For{$i \gets 1$ \To $n$}
        \State $m[i, i] \gets 0$
    \EndFor
    \For{$\ell \gets 2$ \To $n$} \Comment{鏈長}
        \For{$i \gets 1$ \To $n - \ell + 1$}
            \State $j \gets i + \ell - 1$
            \State $m[i, j] \gets \infty$
            \For{$k \gets i$ \To $j - 1$}
                \State $q \gets m[i, k] + m[k + 1, j] + p_{i-1} p_k p_j$
                \If{$q < m[i, j]$}
                    \State $m[i, j] \gets q$
                    \State $s[i, j] \gets k$
                \EndIf
            \EndFor
        \EndFor
    \EndFor
    \State \Return $m$ 和 $s$
\EndFunction
\end{pseudo}

第 10 行用的是嚴格的 $<$：幾個 $k$ 的成本一樣時，$s[i, j]$ 保留最先試到、也就是最小的
那個 $k$。 這是本書的平手規則。

\textbf{時間}：第 9 行對每一組 $(i, j)$ 執行 $j - i = \ell - 1$ 次，所以總次數是
\[
\sum_{\ell=2}^{n} (n - \ell + 1)(\ell - 1) = \frac{n^3 - n}{6} = \Theta(n^3),
\]
空間是 $m$、$s$ 兩張表，$\Theta(n^2)$。 （和式可以令 $t = \ell - 1$ 寫成
$\sum_{t=1}^{n-1} (n - t)t$ 再用 $\sum t$、$\sum t^2$ 的公式算。）

\subsection{手算追蹤}
\label{sec:dp-mc-trace}

用 $\langle 6, 10, 3, 12, 5, 8, 4 \rangle$ 跑一次。 下表一列是一格，寫出每一個 $k$ 的三項
$m[i, k] + m[k+1, j] + p_{i-1}p_k p_j$，粗體是最小的那個，橫線分開不同的鏈長。

{\small
\begin{center}
\begin{tabular}{c p{0.72\linewidth} r c}
\toprule
$[i, j]$ & 每個 $k$ 的成本 & $m$ & $s$ \\
\midrule
$[1,2]$ & $k=1$：$0+0+6\cdot10\cdot3=\mathbf{180}$ & 180 & 1 \\
$[2,3]$ & $k=2$：$0+0+10\cdot3\cdot12=\mathbf{360}$ & 360 & 2 \\
$[3,4]$ & $k=3$：$0+0+3\cdot12\cdot5=\mathbf{180}$ & 180 & 3 \\
$[4,5]$ & $k=4$：$0+0+12\cdot5\cdot8=\mathbf{480}$ & 480 & 4 \\
$[5,6]$ & $k=5$：$0+0+5\cdot8\cdot4=\mathbf{160}$ & 160 & 5 \\
\midrule
$[1,3]$ & $k=1$：$0+360+6\cdot10\cdot12=1080$；$k=2$：$180+0+6\cdot3\cdot12=\mathbf{396}$ & 396 & 2 \\
$[2,4]$ & $k=2$：$0+180+10\cdot3\cdot5=\mathbf{330}$；$k=3$：$360+0+10\cdot12\cdot5=960$ & 330 & 2 \\
$[3,5]$ & $k=3$：$0+480+3\cdot12\cdot8=768$；$k=4$：$180+0+3\cdot5\cdot8=\mathbf{300}$ & 300 & 4 \\
$[4,6]$ & $k=4$：$0+160+12\cdot5\cdot4=\mathbf{400}$；$k=5$：$480+0+12\cdot8\cdot4=864$ & 400 & 4 \\
\midrule
$[1,4]$ & $k=1$：$0+330+6\cdot10\cdot5=630$；$k=2$：$180+180+6\cdot3\cdot5=\mathbf{450}$；$k=3$：$396+0+6\cdot12\cdot5=756$ & 450 & 2 \\
$[2,5]$ & $k=2$：$0+300+10\cdot3\cdot8=\mathbf{540}$；$k=3$：$360+480+10\cdot12\cdot8=1800$；$k=4$：$330+0+10\cdot5\cdot8=730$ & 540 & 2 \\
$[3,6]$ & $k=3$：$0+400+3\cdot12\cdot4=544$；$k=4$：$180+160+3\cdot5\cdot4=400$；$k=5$：$300+0+3\cdot8\cdot4=\mathbf{396}$ & 396 & 5 \\
\midrule
$[1,5]$ & $k=1$：$0+540+6\cdot10\cdot8=1020$；$k=2$：$180+300+6\cdot3\cdot8=\mathbf{624}$；$k=3$：$396+480+6\cdot12\cdot8=1452$；$k=4$：$450+0+6\cdot5\cdot8=690$ & 624 & 2 \\
$[2,6]$ & $k=2$：$0+396+10\cdot3\cdot4=\mathbf{516}$；$k=3$：$360+400+10\cdot12\cdot4=1240$；$k=4$：$330+160+10\cdot5\cdot4=690$；$k=5$：$540+0+10\cdot8\cdot4=860$ & 516 & 2 \\
\midrule
$[1,6]$ & $k=1$：$0+516+6\cdot10\cdot4=756$；$k=2$：$180+396+6\cdot3\cdot4=\mathbf{648}$；$k=3$：$396+400+6\cdot12\cdot4=1084$；$k=4$：$450+160+6\cdot5\cdot4=730$；$k=5$：$624+0+6\cdot8\cdot4=816$ & 648 & 2 \\
\bottomrule
\end{tabular}
\end{center}
}

整理成兩張表，列是 $i$、欄是 $j$，只用到 $i \le j$ 的上三角。 底色深淺交錯的是不同的對角線
（同一條對角線上的格子鏈長相同）；橘框的是等一下還原括號時會用到的格子。

\begin{center}
\begin{minipage}{\linewidth}
\centering
\begin{tikzpicture}[x=1.05cm, y=0.62cm, baseline=(current bounding box.north)]
\node[font=\small] at (3.50,1.15) {$m[i,j]$};
\node[dp-hd] at (1,0.35) {$j=1$};
\node[dp-hd] at (2,0.35) {$2$};
\node[dp-hd] at (3,0.35) {$3$};
\node[dp-hd] at (4,0.35) {$4$};
\node[dp-hd] at (5,0.35) {$5$};
\node[dp-hd] at (6,0.35) {$6$};
\node[dp-hd, anchor=east] at (0.45,-0.5) {$i=1$};
\node[dp-cd0, minimum width=1.05cm] at (1,-0.5) {0};
\node[dp-cpath, minimum width=1.05cm] at (2,-0.5) {180};
\node[dp-cd0, minimum width=1.05cm] at (3,-0.5) {396};
\node[dp-cd1, minimum width=1.05cm] at (4,-0.5) {450};
\node[dp-cd0, minimum width=1.05cm] at (5,-0.5) {624};
\node[dp-cpath, minimum width=1.05cm] at (6,-0.5) {648};
\node[dp-hd, anchor=east] at (0.45,-1.5) {$2$};
\node[dp-cd0, minimum width=1.05cm] at (2,-1.5) {0};
\node[dp-cd1, minimum width=1.05cm] at (3,-1.5) {360};
\node[dp-cd0, minimum width=1.05cm] at (4,-1.5) {330};
\node[dp-cd1, minimum width=1.05cm] at (5,-1.5) {540};
\node[dp-cd0, minimum width=1.05cm] at (6,-1.5) {516};
\node[dp-hd, anchor=east] at (0.45,-2.5) {$3$};
\node[dp-cd0, minimum width=1.05cm] at (3,-2.5) {0};
\node[dp-cpath, minimum width=1.05cm] at (4,-2.5) {180};
\node[dp-cpath, minimum width=1.05cm] at (5,-2.5) {300};
\node[dp-cpath, minimum width=1.05cm] at (6,-2.5) {396};
\node[dp-hd, anchor=east] at (0.45,-3.5) {$4$};
\node[dp-cd0, minimum width=1.05cm] at (4,-3.5) {0};
\node[dp-cd1, minimum width=1.05cm] at (5,-3.5) {480};
\node[dp-cd0, minimum width=1.05cm] at (6,-3.5) {400};
\node[dp-hd, anchor=east] at (0.45,-4.5) {$5$};
\node[dp-cd0, minimum width=1.05cm] at (5,-4.5) {0};
\node[dp-cd1, minimum width=1.05cm] at (6,-4.5) {160};
\node[dp-hd, anchor=east] at (0.45,-5.5) {$6$};
\node[dp-cd0, minimum width=1.05cm] at (6,-5.5) {0};
\end{tikzpicture}\hspace{0.8cm}\begin{tikzpicture}[x=0.8cm, y=0.62cm, baseline=(current bounding box.north)]
\node[font=\small] at (3.00,1.15) {$s[i,j]$};
\node[dp-hd] at (1,0.35) {$j=2$};
\node[dp-hd] at (2,0.35) {$3$};
\node[dp-hd] at (3,0.35) {$4$};
\node[dp-hd] at (4,0.35) {$5$};
\node[dp-hd] at (5,0.35) {$6$};
\node[dp-hd, anchor=east] at (0.45,-0.5) {$i=1$};
\node[dp-cpath, minimum width=0.8cm] at (1,-0.5) {1};
\node[dp-cd0, minimum width=0.8cm] at (2,-0.5) {2};
\node[dp-cd1, minimum width=0.8cm] at (3,-0.5) {2};
\node[dp-cd0, minimum width=0.8cm] at (4,-0.5) {2};
\node[dp-cpath, minimum width=0.8cm] at (5,-0.5) {2};
\node[dp-hd, anchor=east] at (0.45,-1.5) {$2$};
\node[dp-cd1, minimum width=0.8cm] at (2,-1.5) {2};
\node[dp-cd0, minimum width=0.8cm] at (3,-1.5) {2};
\node[dp-cd1, minimum width=0.8cm] at (4,-1.5) {2};
\node[dp-cd0, minimum width=0.8cm] at (5,-1.5) {2};
\node[dp-hd, anchor=east] at (0.45,-2.5) {$3$};
\node[dp-cpath, minimum width=0.8cm] at (3,-2.5) {3};
\node[dp-cpath, minimum width=0.8cm] at (4,-2.5) {4};
\node[dp-cpath, minimum width=0.8cm] at (5,-2.5) {5};
\node[dp-hd, anchor=east] at (0.45,-3.5) {$4$};
\node[dp-cd1, minimum width=0.8cm] at (4,-3.5) {4};
\node[dp-cd0, minimum width=0.8cm] at (5,-3.5) {4};
\node[dp-hd, anchor=east] at (0.45,-4.5) {$5$};
\node[dp-cd1, minimum width=0.8cm] at (5,-4.5) {5};
\end{tikzpicture}
\diagcaption{$\langle 6, 10, 3, 12, 5, 8, 4 \rangle$ 的 $m$ 表和 $s$ 表。 $m[1, 6] = 648$ 是答
案。 橘框是 \proc{Print-Parens} 走過的格子：$[1,6]$、$[1,2]$、$[3,6]$、$[3,5]$、$[3,4]$。}
\label{fig:dp-mc-tables}
\end{minipage}
\end{center}

從表裡還可以看出幾件事：
\begin{itemize}
\item $m[3, 6]$ 的三個候選是 $544$、$400$、$396$，最後一個只贏了 $4$。 鏈長 $4$ 以上的格子，
候選的差距常常很小，手算要特別小心。
\item $s[1, j] = 2$ 對所有 $j \ge 3$ 都成立：$p_2 = 3$ 是整個序列最小的維度，在它這裡切開，
最後一次乘法的 $p_{i-1}p_k p_j$ 就很小。 這只是這組數字的現象，不能拿來當規則：例如
$\langle 2, 3, 4, 5 \rangle$ 最小的中間維度是 $p_1 = 3$，在那裡切開的 $A_1(A_2 A_3)$ 要
$60 + 30 = 90$ 次，$(A_1 A_2) A_3$ 只要 $24 + 40 = 64$ 次。 子問題還沒算完之前，看不出哪個
$k$ 最好，這正是要把每個 $k$ 都試一遍的原因。
\item 鏈長 $2$ 的格子只有一個候選，所以 $s[i, i+1] = i$。
\end{itemize}

\subsection{還原括號}

$s[i, j] = k$ 表示 $A_{i..j}$ 最後一次乘法在 $A_k$ 和 $A_{k+1}$ 之間，所以從 $s[1, n]$ 開
始遞迴地往兩邊拆：

\begin{pseudo}[label=alg:dp-mc-parens]{印出最佳的括號}
\Procedure{Print-Parens}{$s, i, j$}
    \If{$i = j$}
        \State 印出 $A_i$
    \Else
        \State 印出「(」
        \State \Call{Print-Parens}{$s, i, s[i, j]$}
        \State \Call{Print-Parens}{$s, s[i, j] + 1, j$}
        \State 印出「)」
    \EndIf
\EndProcedure
\end{pseudo}

這組數字：$s[1, 6] = 2$，拆成 $A_{1..2}$ 和 $A_{3..6}$；$s[3, 6] = 5$，拆成 $A_{3..5}$ 和
$A_6$；$s[3, 5] = 4$，拆成 $A_{3..4}$ 和 $A_5$。 印出來是
\[
((A_1 A_2)(((A_3 A_4) A_5) A_6)) .
\]

驗算一次：$A_1 A_2$ 花 $6 \cdot 10 \cdot 3 = 180$ 得到 $6 \times 3$；$A_3 A_4$ 花
$3 \cdot 12 \cdot 5 = 180$ 得到 $3 \times 5$；再乘 $A_5$ 花 $3 \cdot 5 \cdot 8 = 120$ 得到
$3 \times 8$；再乘 $A_6$ 花 $3 \cdot 8 \cdot 4 = 96$ 得到 $3 \times 4$；最後 $6 \times 3$ 乘
$3 \times 4$ 花 $72$。 總共 $180 + 180 + 120 + 96 + 72 = 648$，和 $m[1, 6]$ 一樣。

\subsection{C 程式}

\begin{cbookrecap}[第~9~章 9.7 節、第~19~章 19.3 節]
C 的二維陣列要在編譯時知道列寬。 列寬要等執行時才知道時，本書一律用 \texttt{calloc} 要一
塊 $(n+1)(n+1)$ 格的一維陣列，第 $i$ 列第 $j$ 格放在 \texttt{t[i * (n + 1) + j]}，這就是第~9~
章「多維陣列在記憶體裡其實攤平成一列」的寫法。 \texttt{calloc} 會把整塊設成 $0$，用完要
\texttt{free}。
\end{cbookrecap}

$m$、$s$ 兩張表都這樣攤平，用巨集 \texttt{AT(t, i, j)} 取第 $i$ 列第 $j$ 格；索引和虛擬碼一
樣從 $1$ 開始，第 $0$ 列、第 $0$ 欄空著不用，比較好對照。 $m$ 用 $-1$ 表示「還沒有候選」，
取代虛擬碼的 $\infty$。

\begin{example}[label=ex:dp-mc]{c}{(matrix\_chain.c) 矩陣鏈乘的填表與括號}
#include <stdio.h>
#include <stdlib.h>

#define MAXN 100
#define AT(t, i, j) ((t)[(i) * (n + 1) + (j)])   // (n+1) x (n+1) 的表攤平

// 依 s 表印出 A_i..A_j 的最佳括號
void print_parens(const int *s, int n, int i, int j)
{
    if (i == j) {
        printf("A%d", i);
        return;
    }
    printf("(");
    print_parens(s, n, i, AT(s, i, j));
    print_parens(s, n, AT(s, i, j) + 1, j);
    printf(")");
}

int main(void)
{
    int n;
    long long p[MAXN + 1];
    if (scanf("%d", &n) != 1 || n < 1 || n > MAXN) {
        return 1;
    }
    for (int i = 0; i <= n; i++) {
        if (scanf("%lld", &p[i]) != 1) {
            return 1;
        }
    }
    long long *m = calloc((size_t)(n + 1) * (n + 1), sizeof *m);
    int *s = calloc((size_t)(n + 1) * (n + 1), sizeof *s);
    if (m == NULL || s == NULL) {
        free(m);
        free(s);
        return 1;
    }
    for (int len = 2; len <= n; len++) {           // 鏈長由短到長
        for (int i = 1; i + len - 1 <= n; i++) {
            int j = i + len - 1;
            AT(m, i, j) = -1;                      // -1 代表還沒有候選
            for (int k = i; k < j; k++) {
                long long q = AT(m, i, k) + AT(m, k + 1, j)
                            + p[i - 1] * p[k] * p[j];
                if (AT(m, i, j) < 0 || q < AT(m, i, j)) {
                    AT(m, i, j) = q;               // 平手時保留較小的 k
                    AT(s, i, j) = k;
                }
            }
        }
    }
    printf("m[1][%d] = %lld\n", n, AT(m, 1, n));
    print_parens(s, n, 1, n);
    printf("\n");
    free(m);
    free(s);
    return 0;
}
\end{example}
\iosample{6 \\ 6 10 3 12 5 8 4}{m[1][6] = 648 \\ ((A1A2)(((A3A4)A5)A6))}

如果把外面兩層迴圈改成「$i$ 從 $1$ 到 $n$、
$j$ 從 $i + 1$ 到 $n$」一列一列由上往下填，算 $m[1, 3]$ 時 $m[2, 3]$ 還沒算，答案就錯了。
除了依鏈長，另一種正確的順序是 $i$ 從 $n$ \textbf{往下}到 $1$、$j$ 從 $i + 1$ 往上：這樣
算 $m[i, j]$ 時，第 $i + 1$ 列以下都已經填好，第 $i$ 列左邊的格子也填好了。

\section{最長共同子序列}
\label{sec:dp-lcs}

\subsection{定義}

從序列 $X = \langle x_1, x_2, \ldots, x_m \rangle$ 裡刪掉任意幾個元素（可以不刪），剩下的元素保
持原來的順序，得到的序列叫 $X$ 的\textbf{子序列}（subsequence）。 例如 \texttt{ALE} 是
\texttt{PALETTE} 的子序列（取第 $2, 3, 4$ 個），\texttt{LAP} 不是。 子序列不要求連續，這是
它和「子字串」不同的地方。 $Z$ 同時是 $X$ 和 $Y$ 的子序列時，叫\textbf{共同子序列}；最長
的共同子序列叫 \textbf{LCS}（longest common subsequence）。 LCS 不一定只有一個，但長度是
確定的。

章首的問題就是 LCS：把每一行程式當成一個元素，兩個版本的 LCS 就是最多能原封不動對上的
行。 Unix 的 \texttt{diff}、版本控制系統比較兩個檔案，背後都是這個問題（實際的工具用了更
快的變形）。 生物資訊裡比對兩段 DNA，也常用 LCS 和下一節的編輯距離。

暴力法是列出 $X$ 的 $2^m$ 個子序列，一個一個檢查是不是 $Y$ 的子序列（從左到右掃一次
$Y$，$O(n)$），時間 $\Theta(n 2^m)$。 下面用\textbf{前綴}當子問題：$X_i = \langle x_1,
\ldots, x_i \rangle$ 是 $X$ 的前 $i$ 個元素，$X_0$ 是空序列。

\subsection{最後一個字元的三種情形}

$Z$ 是 $X$ 和 $Y$ 的共同子序列，表示 $Z$ 的每個元素可以依序\textbf{配}到 $X$ 的一個位置和
$Y$ 的一個位置，配到的位置在 $X$、$Y$ 裡都是由左往右遞增的。 只有 $Z$ 的最後一個元素能配到
$x_m$，也只有它能配到 $y_n$。 所以對任何一種配法，下面三種情形一定有一種成立：

\begin{enumerate}
\item $x_m$ 和 $y_n$ 互相配對（這時一定 $x_m = y_n$）；
\item $x_m$ 沒有被配到；
\item $y_n$ 沒有被配到。
\end{enumerate}

令 $c[i, j]$ 為 $X_i$ 和 $Y_j$ 的 LCS 長度。 由這三種情形得到：

\begin{theorem}
\label{thm:dp-lcs}
$i, j \ge 1$ 時，
\[
c[i, j] = \begin{cases}
c[i - 1, j - 1] + 1, & x_i = y_j, \\
\max\bigl(c[i - 1, j],\ c[i, j - 1]\bigr), & x_i \ne y_j,
\end{cases}
\]
而 $c[i, 0] = c[0, j] = 0$。
\end{theorem}

\begin{proof}
先證 $x_i \ne y_j$ 的情形。 $X_{i-1}$ 和 $Y_j$ 的任何共同子序列也是 $X_i$ 和 $Y_j$ 的，所以
$c[i, j] \ge c[i - 1, j]$；同理 $c[i, j] \ge c[i, j - 1]$。 反過來，取 $X_i$、$Y_j$ 的一個
LCS 和它的一種配法。 情形 1 不可能（$x_i \ne y_j$），所以 $x_i$ 或 $y_j$ 沒有被配到；前者
表示它是 $X_{i-1}$ 和 $Y_j$ 的共同子序列，長度不超過 $c[i-1, j]$，後者長度不超過
$c[i, j-1]$。 兩個方向合起來就是等式。

再證 $x_i = y_j$。 在 $X_{i-1}$、$Y_{j-1}$ 的 LCS 後面接上 $x_i$，得到 $X_i$、$Y_j$ 的共同
子序列，所以 $c[i, j] \ge c[i-1, j-1] + 1$。 反過來，取一個 LCS 和它的配法：情形 1 時去掉
最後一個元素，剩下的是 $X_{i-1}$、$Y_{j-1}$ 的共同子序列，長度 $\le c[i-1, j-1]$；情形 2
時它是 $X_{i-1}$、$Y_j$ 的共同子序列，而 $X_{i-1}$ 和 $Y_j$ 的共同子序列去掉最後一個元素
（如果它配到 $y_j$）就是 $X_{i-1}$、$Y_{j-1}$ 的共同子序列，所以長度
$\le c[i-1, j-1] + 1$；情形 3 對稱。 三種情形都不超過 $c[i-1, j-1] + 1$。
\end{proof}

這個定理也說明了最佳子結構：LCS 的前面一段，一定是某一對前綴的 LCS。 子問題有
$(m+1)(n+1)$ 個，每個只看三個鄰居，所以填表是 $\Theta(mn)$。 $x_i = y_j$ 時\textbf{不用}
再和 $c[i-1, j]$、$c[i, j-1]$ 比，配對一定不會比較差，這是證明裡的第二段。

\subsection{填表與箭頭}

除了數字，每一格再記一個箭頭 $b[i, j]$，說明這一格的值是從哪裡來的：$\nwarrow$ 表示
$x_i = y_j$、從左上角加一；$\uparrow$ 表示取上面的 $c[i-1, j]$；$\leftarrow$ 表示取左邊的
$c[i, j-1]$。

\begin{pseudo}[label=alg:dp-lcs-length]{LCS 的填表}
\Function{LCS-Length}{$X, Y, m, n$}
    \For{$i \gets 0$ \To $m$}
        \State $c[i, 0] \gets 0$
    \EndFor
    \For{$j \gets 0$ \To $n$}
        \State $c[0, j] \gets 0$
    \EndFor
    \For{$i \gets 1$ \To $m$}
        \For{$j \gets 1$ \To $n$}
            \If{$x_i = y_j$}
                \State $c[i, j] \gets c[i - 1, j - 1] + 1$
                \State $b[i, j] \gets \text{「}\nwarrow\text{」}$
            \ElsIf{$c[i - 1, j] \ge c[i, j - 1]$}
                \State $c[i, j] \gets c[i - 1, j]$
                \State $b[i, j] \gets \text{「}\uparrow\text{」}$
            \Else
                \State $c[i, j] \gets c[i, j - 1]$
                \State $b[i, j] \gets \text{「}\leftarrow\text{」}$
            \EndIf
        \EndFor
    \EndFor
    \State \Return $c$ 和 $b$
\EndFunction
\end{pseudo}

第 11 行用 $\ge$：上面和左邊一樣大時，箭頭指向上。 這是本書（也是 CLRS）的平手規則。

\begin{pseudo}[label=alg:dp-print-lcs]{沿箭頭印出一個 LCS}
\Procedure{Print-LCS}{$b, X, i, j$}
    \If{$i = 0$ \OR $j = 0$}
        \State \Return
    \EndIf
    \If{$b[i, j] = \text{「}\nwarrow\text{」}$}
        \State \Call{Print-LCS}{$b, X, i - 1, j - 1$}
        \State 印出 $x_i$
    \ElsIf{$b[i, j] = \text{「}\uparrow\text{」}$}
        \State \Call{Print-LCS}{$b, X, i - 1, j$}
    \Else
        \State \Call{Print-LCS}{$b, X, i, j - 1$}
    \EndIf
\EndProcedure
\end{pseudo}

呼叫 $\proc{Print-LCS}(b, X, m, n)$。 每次呼叫 $i$ 或 $j$ 至少少一，所以是 $O(m + n)$。 先遞
迴再印，印出來的順序才是由左到右。

\subsection{手算追蹤}
\label{sec:dp-lcs-trace}

$X = \texttt{PALETTE}$（$m = 7$），$Y = \texttt{PLANET}$（$n = 6$）。 列是 $X$、欄是 $Y$，
一列一列由左往右填。

\begin{center}
\begin{minipage}{\linewidth}
\centering
\begin{tikzpicture}[x=0.98cm, y=0.66cm]
\node[dp-hd] at (0,1) {$j$};
\node[dp-hd] at (-1,0) {$i$};
\node[dp-hd] at (1,1) {0};
\node[dp-hd] at (2,1) {1};
\node[dp-hd] at (2,0.4) {\texttt{P}};
\node[dp-hd] at (3,1) {2};
\node[dp-hd] at (3,0.4) {\texttt{L}};
\node[dp-hd] at (4,1) {3};
\node[dp-hd] at (4,0.4) {\texttt{A}};
\node[dp-hd] at (5,1) {4};
\node[dp-hd] at (5,0.4) {\texttt{N}};
\node[dp-hd] at (6,1) {5};
\node[dp-hd] at (6,0.4) {\texttt{E}};
\node[dp-hd] at (7,1) {6};
\node[dp-hd] at (7,0.4) {\texttt{T}};
\node[dp-hd] at (-1,-0.6) {0};
\node[dp-c0, minimum width=0.98cm] at (1,-0.6) {0};
\node[dp-c0, minimum width=0.98cm] at (2,-0.6) {0};
\node[dp-c0, minimum width=0.98cm] at (3,-0.6) {0};
\node[dp-c0, minimum width=0.98cm] at (4,-0.6) {0};
\node[dp-c0, minimum width=0.98cm] at (5,-0.6) {0};
\node[dp-c0, minimum width=0.98cm] at (6,-0.6) {0};
\node[dp-c0, minimum width=0.98cm] at (7,-0.6) {0};
\node[dp-hd] at (-1,-1.6) {1};
\node[dp-hd] at (-0.3,-1.6) {\texttt{P}};
\node[dp-c0, minimum width=0.98cm] at (1,-1.6) {0};
\node[dp-chit, minimum width=0.98cm] at (2,-1.6) {$\nwarrow$\,1};
\node[dp-cpath, minimum width=0.98cm] at (3,-1.6) {$\leftarrow$\,1};
\node[dp-c, minimum width=0.98cm] at (4,-1.6) {$\leftarrow$\,1};
\node[dp-c, minimum width=0.98cm] at (5,-1.6) {$\leftarrow$\,1};
\node[dp-c, minimum width=0.98cm] at (6,-1.6) {$\leftarrow$\,1};
\node[dp-c, minimum width=0.98cm] at (7,-1.6) {$\leftarrow$\,1};
\node[dp-hd] at (-1,-2.6) {2};
\node[dp-hd] at (-0.3,-2.6) {\texttt{A}};
\node[dp-c0, minimum width=0.98cm] at (1,-2.6) {0};
\node[dp-c, minimum width=0.98cm] at (2,-2.6) {$\uparrow$\,1};
\node[dp-c, minimum width=0.98cm] at (3,-2.6) {$\uparrow$\,1};
\node[dp-chit, minimum width=0.98cm] at (4,-2.6) {$\nwarrow$\,2};
\node[dp-cpath, minimum width=0.98cm] at (5,-2.6) {$\leftarrow$\,2};
\node[dp-c, minimum width=0.98cm] at (6,-2.6) {$\leftarrow$\,2};
\node[dp-c, minimum width=0.98cm] at (7,-2.6) {$\leftarrow$\,2};
\node[dp-hd] at (-1,-3.6) {3};
\node[dp-hd] at (-0.3,-3.6) {\texttt{L}};
\node[dp-c0, minimum width=0.98cm] at (1,-3.6) {0};
\node[dp-c, minimum width=0.98cm] at (2,-3.6) {$\uparrow$\,1};
\node[dp-c, minimum width=0.98cm] at (3,-3.6) {$\nwarrow$\,2};
\node[dp-c, minimum width=0.98cm] at (4,-3.6) {$\uparrow$\,2};
\node[dp-cpath, minimum width=0.98cm] at (5,-3.6) {$\uparrow$\,2};
\node[dp-c, minimum width=0.98cm] at (6,-3.6) {$\uparrow$\,2};
\node[dp-c, minimum width=0.98cm] at (7,-3.6) {$\uparrow$\,2};
\node[dp-hd] at (-1,-4.6) {4};
\node[dp-hd] at (-0.3,-4.6) {\texttt{E}};
\node[dp-c0, minimum width=0.98cm] at (1,-4.6) {0};
\node[dp-c, minimum width=0.98cm] at (2,-4.6) {$\uparrow$\,1};
\node[dp-c, minimum width=0.98cm] at (3,-4.6) {$\uparrow$\,2};
\node[dp-c, minimum width=0.98cm] at (4,-4.6) {$\uparrow$\,2};
\node[dp-c, minimum width=0.98cm] at (5,-4.6) {$\uparrow$\,2};
\node[dp-chit, minimum width=0.98cm] at (6,-4.6) {$\nwarrow$\,3};
\node[dp-c, minimum width=0.98cm] at (7,-4.6) {$\leftarrow$\,3};
\node[dp-hd] at (-1,-5.6) {5};
\node[dp-hd] at (-0.3,-5.6) {\texttt{T}};
\node[dp-c0, minimum width=0.98cm] at (1,-5.6) {0};
\node[dp-c, minimum width=0.98cm] at (2,-5.6) {$\uparrow$\,1};
\node[dp-c, minimum width=0.98cm] at (3,-5.6) {$\uparrow$\,2};
\node[dp-c, minimum width=0.98cm] at (4,-5.6) {$\uparrow$\,2};
\node[dp-c, minimum width=0.98cm] at (5,-5.6) {$\uparrow$\,2};
\node[dp-cpath, minimum width=0.98cm] at (6,-5.6) {$\uparrow$\,3};
\node[dp-c, minimum width=0.98cm] at (7,-5.6) {$\nwarrow$\,4};
\node[dp-hd] at (-1,-6.6) {6};
\node[dp-hd] at (-0.3,-6.6) {\texttt{T}};
\node[dp-c0, minimum width=0.98cm] at (1,-6.6) {0};
\node[dp-c, minimum width=0.98cm] at (2,-6.6) {$\uparrow$\,1};
\node[dp-c, minimum width=0.98cm] at (3,-6.6) {$\uparrow$\,2};
\node[dp-c, minimum width=0.98cm] at (4,-6.6) {$\uparrow$\,2};
\node[dp-c, minimum width=0.98cm] at (5,-6.6) {$\uparrow$\,2};
\node[dp-c, minimum width=0.98cm] at (6,-6.6) {$\uparrow$\,3};
\node[dp-chit, minimum width=0.98cm] at (7,-6.6) {$\nwarrow$\,4};
\node[dp-hd] at (-1,-7.6) {7};
\node[dp-hd] at (-0.3,-7.6) {\texttt{E}};
\node[dp-c0, minimum width=0.98cm] at (1,-7.6) {0};
\node[dp-c, minimum width=0.98cm] at (2,-7.6) {$\uparrow$\,1};
\node[dp-c, minimum width=0.98cm] at (3,-7.6) {$\uparrow$\,2};
\node[dp-c, minimum width=0.98cm] at (4,-7.6) {$\uparrow$\,2};
\node[dp-c, minimum width=0.98cm] at (5,-7.6) {$\uparrow$\,2};
\node[dp-c, minimum width=0.98cm] at (6,-7.6) {$\nwarrow$\,3};
\node[dp-cpath, minimum width=0.98cm] at (7,-7.6) {$\uparrow$\,4};
\end{tikzpicture}
\diagcaption{\texttt{PALETTE} 和 \texttt{PLANET} 的 $c$ 表和箭頭。 第 $0$ 列、第 $0$ 欄是
邊界。 從右下角 $c[7, 6] = 4$ 沿箭頭走回左上，經過的格子加框；粉紅色的是 $\nwarrow$，也就
是被選進 LCS 的字元：\texttt{P}、\texttt{A}、\texttt{E}、\texttt{T}。}
\label{fig:dp-lcs-table}
\end{minipage}
\end{center}

看幾格是怎麼來的：
\begin{itemize}
\item $c[2, 3]$：$x_2 = \texttt{A} = y_3$，從左上 $c[1, 2] = 1$ 加一，得 $2$，箭頭 $\nwarrow$。
\item $c[3, 3]$：$x_3 = \texttt{L} \ne \texttt{A}$，上面 $c[2, 3] = 2$、左邊 $c[3, 2] = 2$，平
手，照規則取上面，箭頭 $\uparrow$。
\item $c[4, 5]$：$x_4 = \texttt{E} = y_5$，從 $c[3, 4] = 2$ 加一，得 $3$。
\item $c[5, 6]$：$x_5 = \texttt{T} = y_6$，從 $c[4, 5] = 3$ 加一，得 $4$。
\item $c[6, 6]$：$x_6$ 也是 \texttt{T}，從 $c[5, 5] = 3$ 加一，也是 $4$，箭頭 $\nwarrow$。
\item $c[7, 6]$：$x_7 = \texttt{E} \ne \texttt{T}$，上面是 $4$、左邊是 $3$，取上面。
\end{itemize}

從 $c[7, 6]$ 往回走：$\uparrow$ 到 $[6, 6]$，$\nwarrow$（選 $x_6 = \texttt{T}$）到 $[5, 5]$，
$\uparrow$ 到 $[4, 5]$，$\nwarrow$（選 \texttt{E}）到 $[3, 4]$，$\uparrow$ 到 $[2, 4]$，
$\leftarrow$ 到 $[2, 3]$，$\nwarrow$（選 \texttt{A}）到 $[1, 2]$，$\leftarrow$ 到
$[1, 1]$，$\nwarrow$（選 \texttt{P}）到 $[0, 0]$。 倒過來讀，LCS 是 \texttt{PAET}。

這兩個字串還有另一個長度 $4$ 的 LCS：\texttt{PLET}（\texttt{PALETTE} 的第 $1, 3, 4, 5$ 個字
元）。 沿著箭頭只能走出一條路，所以 \proc{Print-LCS} 只印出其中一個；印出哪一個由平手規
則決定。

\begin{examnote}
平手時選哪一邊，不同的課本規定不同。 LCS 平手時本書指向上（第 11 行的 $\ge$），有些書指
向左；同一題 \texttt{PALETTE}、\texttt{PLANET}，指向左會還原出 \texttt{PLET}。 矩陣鏈乘
和最佳二元搜尋樹，本書在平手時取最小的 $k$ 或 $r$。 考題有規定時照題目；沒有規定時，在答
案裡寫一句你用的規則（例如「平手時取 $\uparrow$」），長度和最佳值不受影響，只有還原出來
的解可能不同。 題目問「所有的 LCS」時，平手的格子兩個方向都要走。
\end{examnote}

\subsection{省掉箭頭表}

$b$ 表其實可以不存。 還原時站在 $[i, j]$，只要重新判斷一次：$x_i = y_j$ 就往左上；否則比較
$c[i-1, j]$ 和 $c[i, j-1]$，用和填表時一樣的平手規則決定往上還是往左。 每一步 $O(1)$，結
果和用箭頭完全相同，省下一張 $(m+1)(n+1)$ 的表。

\subsection{只留兩列}
\label{sec:dp-lcs-space}

如果只要 LCS 的長度，連 $c$ 表都不用整張留著。 第 $i$ 列只用到第 $i - 1$ 列和自己左邊的格
子，所以只要兩列：$\mathit{prev}$ 是上一列，$\mathit{cur}$ 是正在填的這一列，一列填完後交換
兩者的角色。

\begin{pseudo}[label=alg:dp-lcs-two-rows]{只用兩列算 LCS 長度}
\Function{LCS-Length-Two-Rows}{$X, Y, m, n$}
    \State 建立 $\mathit{prev}[0..n]$ 和 $\mathit{cur}[0..n]$，每一格都設成 $0$
    \For{$i \gets 1$ \To $m$}
        \For{$j \gets 1$ \To $n$}
            \If{$x_i = y_j$}
                \State $\mathit{cur}[j] \gets \mathit{prev}[j - 1] + 1$
            \Else
                \State $\mathit{cur}[j] \gets \max(\mathit{prev}[j], \mathit{cur}[j - 1])$
            \EndIf
        \EndFor
        \State 交換 $\mathit{prev}$ 和 $\mathit{cur}$ \Comment{只交換指標，不複製}
    \EndFor
    \State \Return $\mathit{prev}[n]$
\EndFunction
\end{pseudo}

$\mathit{cur}[0]$ 從來沒被寫過，一直是 $0$，正好是邊界 $c[i, 0]$。 第 9 行換完之後，
$\mathit{prev}$ 存的是剛填好的第 $i$ 列，所以迴圈結束時答案在 $\mathit{prev}[n]$。 空間從
$\Theta(mn)$ 降到 $\Theta(n)$；讓較短的字串當 $Y$，就是 $\Theta(\min(m, n))$。 再省一點，只用一
列也可以：由左往右更新 $\mathit{row}[j]$ 之前，先把舊的 $\mathit{row}[j]$ 存進一個變數，填下一
格時它就是需要的左上角 $c[i-1, j]$。

代價是\textbf{還原不出 LCS 本身}：走回去時需要的那些格子已經被蓋掉了。 要在線性空間裡同時
得到 LCS，可以用 Hirschberg 的方法：用兩列的版本從前面、從後面各算一次，找出 LCS 通過第
$\lfloor m/2 \rfloor$ 列的哪一格，再對左上、右下兩塊分治。 時間仍是 $\Theta(mn)$，這裡不
細講。

\subsection{C 程式}

\texttt{lcs} 填整張表，照前面「省掉箭頭表」的方法直接用 $c$ 表還原。 只要長度的話，把表
換成虛擬碼~\ref{alg:dp-lcs-two-rows} 的兩列即可。 字串的第 $i$ 個字元在 C 裡是 \texttt{x[i - 1]}，表的索引仍然從 $0$ 開始，第 $0$
列、第 $0$ 欄是邊界。

\begin{example}[label=ex:dp-lcs]{c}{(lcs.c) LCS 的填表與還原}
#include <stdio.h>
#include <stdlib.h>
#include <string.h>

#define MAXLEN 1000

// 填完整的 (m+1) x (n+1) 表，再從右下角走回來還原一個 LCS 放進 out。
int lcs(const char *x, const char *y, char *out)
{
    int m = strlen(x), n = strlen(y);
    int *c = calloc((size_t)(m + 1) * (n + 1), sizeof *c);
    if (c == NULL) {
        return -1;
    }
#define C(i, j) c[(i) * (n + 1) + (j)]
    for (int i = 1; i <= m; i++) {
        for (int j = 1; j <= n; j++) {
            if (x[i - 1] == y[j - 1]) {
                C(i, j) = C(i - 1, j - 1) + 1;
            } else if (C(i - 1, j) >= C(i, j - 1)) {
                C(i, j) = C(i - 1, j);          // 平手時往上
            } else {
                C(i, j) = C(i, j - 1);
            }
        }
    }
    int len = C(m, n);
    int i = m, j = n, k = len;
    out[k] = '\0';
    while (i > 0 && j > 0) {                    // 照填表時的規則走回去
        if (x[i - 1] == y[j - 1]) {
            out[--k] = x[i - 1];
            i--;
            j--;
        } else if (C(i - 1, j) >= C(i, j - 1)) {
            i--;
        } else {
            j--;
        }
    }
#undef C
    free(c);
    return len;
}

int main(void)
{
    static char x[MAXLEN + 1], y[MAXLEN + 1], z[MAXLEN + 1];
    while (scanf("%1000s %1000s", x, y) == 2) {
        int len = lcs(x, y, z);
        printf("length %d: %s\n", len, z);
    }
    return 0;
}
\end{example}
\iosample{PALETTE PLANET \\ PLANET PALETTE \\ SUBSEQUENCE CONSEQUENCES}{length 4: PAET \\ length 4: PLET \\ length 8: SEQUENCE}

第二組把兩個字串對調，長度一樣是 $4$，還原出來的卻是 \texttt{PLET}：對調之後列和欄互換，
「平手取上」變成了原本的「平手取左」。 第三組說明 LCS 不必是連續的一段：
\texttt{SUBSEQUENCE} 和 \texttt{CONSEQUENCES} 的 LCS 是 \texttt{SEQUENCE}，長度 $8$。

\section{最佳二元搜尋樹}
\label{sec:dp-obst}

\subsection{問題與成本模型}

編譯器讀程式時，每讀到一個詞，就要判斷它是不是關鍵字。 假設只看五個關鍵字，依字典順序是
\[
k_1 = \texttt{char} < k_2 = \texttt{for} < k_3 = \texttt{goto} < k_4 = \texttt{int}
< k_5 = \texttt{while},
\]
存在一棵二元搜尋樹裡（第~9~章）。 從大量程式統計出來，每 $50$ 次查詢裡，查到各個關鍵字的
次數 $p_i$，以及查不到、落在兩個關鍵字之間的次數 $q_i$ 如下：

\begin{center}
\begin{tabular}{c*{6}{c}}
\toprule
$i$ & 0 & 1 & 2 & 3 & 4 & 5 \\
\midrule
$k_i$ & & \texttt{char} & \texttt{for} & \texttt{goto} & \texttt{int} & \texttt{while} \\
$p_i$ & & 4 & 10 & 3 & 8 & 6 \\
$q_i$ & 3 & 2 & 4 & 1 & 3 & 6 \\
\bottomrule
\end{tabular}
\end{center}

$q_0$ 是比 $k_1$ 小的詞（例如 \texttt{buf}），$q_i$（$1 \le i < n$）是介於 $k_i$ 和
$k_{i+1}$ 之間的詞（例如 $q_3$ 數的是 \texttt{goto} 和 \texttt{int} 之間的 \texttt{i}、
\texttt{idx}），$q_5$ 是比 \texttt{while} 大的詞（例如 \texttt{x}）。 $\sum p_i + \sum q_i
= 31 + 19 = 50$。 有些書用機率，那只是把每個數除以總數，下面的計算都一樣。

把每個查不到的範圍畫成一個\textbf{外部節點} $d_0, d_1, \ldots, d_n$，接在原本是 \NIL 的位置
上，得到第~7~章的延伸二元樹：$n$ 個鍵是內部節點，$n + 1$ 個外部節點是葉節點，中序走訪的
順序是 $d_0, k_1, d_1, k_2, \ldots, k_n, d_n$。 一次查不到的搜尋，一定停在某個 $d_i$。

\textbf{成本模型。} 本書數的是\textbf{和鍵比較的次數}，深度和高度沿用第~7~章的定義（深度數
邊、根的深度是 $0$；高度數節點）：
\begin{itemize}
\item 搜尋 $k_i$，從根往下走到 $k_i$，路上每個鍵都比一次，所以是 $\operatorname{depth}(k_i)
+ 1$ 次比較，也就是 $k_i$ 所在的\textbf{層數}。
\item 查不到、停在 $d_i$，路上比較的是 $d_i$ 上面的每一個鍵，共 $\operatorname{depth}(d_i)$
次；$d_i$ 本身不是鍵，走到它就知道查不到，不再比較。
\end{itemize}
所以一棵樹 $T$ 的總成本是
\[
\operatorname{cost}(T) = \sum_{i=1}^{n} p_i \bigl(\operatorname{depth}_T(k_i) + 1\bigr)
+ \sum_{i=0}^{n} q_i \operatorname{depth}_T(d_i) ,
\]
除以總次數就是平均每次查詢比較幾次。 查到時最多比較的次數等於內部節點組成的那棵樹的
高度，這正是第~7~章「高度數節點」的好處。

\textbf{最佳二元搜尋樹}（optimal binary search tree）是在所有以 $k_1, \ldots, k_n$ 為鍵的
二元搜尋樹中，總成本最小的一棵。 $n$ 個鍵的二元搜尋樹形狀有卡特蘭數 $C_n$ 種（第~7~章），
全部試一遍是指數時間。

\begin{center}
\begin{minipage}{\linewidth}
\centering
\begin{tikzpicture}[x=0.62cm, y=0.95cm, baseline=(current bounding box.north)]
\node[dp-ext] (gd0) at (0,-2) {$d_0$};
\node[dp-key] (gk1) at (1,-1) {\texttt{char}};
\node[dp-ext] (gd1) at (2,-2) {$d_1$};
\node[dp-key] (gk2) at (3,0) {\texttt{for}};
\node[dp-ext] (gd2) at (4,-3) {$d_2$};
\node[dp-key] (gk3) at (5,-2) {\texttt{goto}};
\node[dp-ext] (gd3) at (6,-3) {$d_3$};
\node[dp-key] (gk4) at (7,-1) {\texttt{int}};
\node[dp-ext] (gd4) at (8,-3) {$d_4$};
\node[dp-key] (gk5) at (9,-2) {\texttt{while}};
\node[dp-ext] (gd5) at (10,-3) {$d_5$};
\draw[dp-e] (gk1) -- (gd0);
\draw[dp-e] (gk2) -- (gk1);
\draw[dp-e] (gk1) -- (gd1);
\draw[dp-e] (gk3) -- (gd2);
\draw[dp-e] (gk4) -- (gk3);
\draw[dp-e] (gk3) -- (gd3);
\draw[dp-e] (gk2) -- (gk4);
\draw[dp-e] (gk5) -- (gd4);
\draw[dp-e] (gk4) -- (gk5);
\draw[dp-e] (gk5) -- (gd5);
\node[font=\small] at (5,-3.8) {(a) 每次挑次數最多的當根，成本 $113$};
\end{tikzpicture}\hspace{0.5cm}\begin{tikzpicture}[x=0.62cm, y=0.95cm, baseline=(current bounding box.north)]
\node[dp-ext] (od0) at (0,-3) {$d_0$};
\node[dp-key] (ok1) at (1,-2) {\texttt{char}};
\node[dp-ext] (od1) at (2,-3) {$d_1$};
\node[dp-key] (ok2) at (3,-1) {\texttt{for}};
\node[dp-ext] (od2) at (4,-3) {$d_2$};
\node[dp-key] (ok3) at (5,-2) {\texttt{goto}};
\node[dp-ext] (od3) at (6,-3) {$d_3$};
\node[dp-key] (ok4) at (7,0) {\texttt{int}};
\node[dp-ext] (od4) at (8,-2) {$d_4$};
\node[dp-key] (ok5) at (9,-1) {\texttt{while}};
\node[dp-ext] (od5) at (10,-2) {$d_5$};
\draw[dp-e] (ok1) -- (od0);
\draw[dp-e] (ok2) -- (ok1);
\draw[dp-e] (ok1) -- (od1);
\draw[dp-e] (ok4) -- (ok2);
\draw[dp-e] (ok3) -- (od2);
\draw[dp-e] (ok2) -- (ok3);
\draw[dp-e] (ok3) -- (od3);
\draw[dp-e] (ok5) -- (od4);
\draw[dp-e] (ok4) -- (ok5);
\draw[dp-e] (ok5) -- (od5);
\node[font=\small] at (5,-3.8) {(b) 最佳二元搜尋樹，成本 $109$};
\end{tikzpicture}
\diagcaption{同一組次數的兩棵樹，圓角框是鍵，灰色方塊是外部節點。 (a) 每次挑次數最多的鍵當
根：先挑 \texttt{for}（$10$ 次），右邊再挑 \texttt{int}。 (b) 最佳的樹，根是
\texttt{int}，只出現 $8$ 次。}
\label{fig:dp-obst-trees}
\end{minipage}
\end{center}

直覺的做法是讓最常查的鍵當根。 圖~\ref{fig:dp-obst-trees}(a) 這樣做，成本 $113$；(b) 的
根是次數第二多的 \texttt{int}，成本卻只有 $109$。 和 (b) 相比，(a) 讓 \texttt{for}、\texttt{char}、$d_0$、$d_1$ 各少比一次，省下
$10 + 4 + 3 + 2 = 19$；卻讓 \texttt{int}、\texttt{while}、$d_4$、$d_5$ 各多比一次，多了
$8 + 6 + 3 + 6 = 23$。 下表逐項算 (b) 的成本，
比較次數一欄對鍵是深度加一、對外部節點是深度：

\begin{center}
\begin{tabular}{l*{11}{c}c}
\toprule
節點 & $d_0$ & $k_1$ & $d_1$ & $k_2$ & $d_2$ & $k_3$ & $d_3$ & $k_4$ & $d_4$ & $k_5$ & $d_5$ & 合計 \\
\midrule
深度 & 3 & 2 & 3 & 1 & 3 & 2 & 3 & 0 & 2 & 1 & 2 & \\
比較次數 & 3 & 3 & 3 & 2 & 3 & 3 & 3 & 1 & 2 & 2 & 2 & \\
次數 & 3 & 4 & 2 & 10 & 4 & 3 & 1 & 8 & 3 & 6 & 6 & 50 \\
乘積 & 9 & 12 & 6 & 20 & 12 & 9 & 3 & 8 & 6 & 12 & 12 & 109 \\
\bottomrule
\end{tabular}
\end{center}

平均每次查詢比 $109 / 50 = 2.18$ 次。 這組次數下，$42$ 種形狀裡沒有比 $109$ 更小的，下面
用動態規劃算出這件事。

\subsection{最佳子結構與遞迴式}

二元搜尋樹的任何一棵子樹，裡面的鍵一定是連續的一段 $k_i, \ldots, k_j$，外部節點也一定是對
應的 $d_{i-1}, \ldots, d_j$（中序順序決定的）。 如果最佳的樹 $T$ 有一棵子樹 $T'$ 裝的是
$k_i..k_j$，那麼 $T'$ 一定是這一段的最佳樹：把 $T'$ 換成這一段更好的樹，$T$ 的其他部分不
變，總成本會下降，矛盾。 這又是剪下貼上。

接著看把一棵子樹掛到一個新的根底下會怎樣：子樹裡每個節點的深度都加一，總成本增加的量正好
是子樹裡所有次數的和。 定義
\[
w(i, j) = \sum_{l=i}^{j} p_l + \sum_{l=i-1}^{j} q_l ,
\]
是裝著 $k_i..k_j$ 的子樹的\textbf{總重量}（注意外部節點從 $d_{i-1}$ 數起）。 令 $e[i, j]$
為 $k_i..k_j$ 的最佳樹的成本。 選 $k_r$（$i \le r \le j$）當根時，左子樹是 $k_i..k_{r-1}$，
右子樹是 $k_{r+1}..k_j$，掛上去之後兩邊各多 $w(i, r-1)$ 和 $w(r+1, j)$，根本身比一次，貢獻
$p_r$。 三者加起來剛好是 $w(i, j)$，所以
\[
e[i, j] = \begin{cases}
0, & j = i - 1, \\[2pt]
\displaystyle\min_{i \le r \le j} \bigl\{\, e[i, r-1] + e[r+1, j] \,\bigr\} + w(i, j), & i \le j .
\end{cases}
\]
$j = i - 1$ 是沒有鍵的「空樹」，只有一個外部節點 $d_{i-1}$，走到它不用比較，成本是 $0$。 另
外用 $\mathit{root}[i, j]$ 記下取到最小值的 $r$。 $w$ 也用表來存，避免每次重加：
\[
w[i, i - 1] = q_{i-1}, \qquad w[i, j] = w[i, j - 1] + p_j + q_j .
\]

\begin{pseudo}[label=alg:dp-obst]{最佳二元搜尋樹}
\Function{Optimal-BST}{$p, q, n$}
    \For{$i \gets 1$ \To $n + 1$}
        \State $e[i, i - 1] \gets 0$
        \State $w[i, i - 1] \gets q_{i-1}$
    \EndFor
    \For{$\ell \gets 1$ \To $n$} \Comment{這一段有幾個鍵}
        \For{$i \gets 1$ \To $n - \ell + 1$}
            \State $j \gets i + \ell - 1$
            \State $e[i, j] \gets \infty$
            \State $w[i, j] \gets w[i, j - 1] + p_j + q_j$
            \For{$r \gets i$ \To $j$}
                \State $t \gets e[i, r - 1] + e[r + 1, j] + w[i, j]$
                \If{$t < e[i, j]$}
                    \State $e[i, j] \gets t$
                    \State $\mathit{root}[i, j] \gets r$
                \EndIf
            \EndFor
        \EndFor
    \EndFor
    \State \Return $e$ 和 $\mathit{root}$
\EndFunction
\end{pseudo}

表的大小：$e$ 和 $w$ 的列 $i$ 從 $1$ 到 $n + 1$、欄 $j$ 從 $0$ 到 $n$（要放得下 $e[n+1, n]$
這種空樹），$\mathit{root}$ 只用 $1 \le i \le j \le n$。 結構和矩陣鏈乘幾乎一樣：依區間長度
填，第 10 行的迴圈每格跑 $\ell$ 次，總時間 $\Theta(n^3)$，空間 $\Theta(n^2)$。 第 12 行用
$<$，平手時保留最小的 $r$。

\subsection{手算追蹤}
\label{sec:dp-obst-trace}

$\ell = 1$ 時只有一個鍵，左右都是空樹，$e[i, i] = w[i, i] = q_{i-1} + p_i + q_i$：
$e[1,1] = 3 + 4 + 2 = 9$、$e[2,2] = 16$、$e[3,3] = 8$、$e[4,4] = 12$、$e[5,5] = 15$，根都是
自己。 $\ell \ge 2$ 的每一格寫出每個 $r$ 的 $e[i, r-1] + e[r+1, j] + w[i, j]$：

{\small
\begin{center}
\begin{tabular}{c c p{0.62\linewidth} c c}
\toprule
$(i, j)$ & $w$ & 每個 $r$ 的成本 & $e$ & $\mathit{root}$ \\
\midrule
$(1,2)$ & 23 & $r=1$：$0+16+23=39$；$r=2$：$9+0+23=\mathbf{32}$ & 32 & 2 \\
$(2,3)$ & 20 & $r=2$：$0+8+20=\mathbf{28}$；$r=3$：$16+0+20=36$ & 28 & 2 \\
$(3,4)$ & 19 & $r=3$：$0+12+19=31$；$r=4$：$8+0+19=\mathbf{27}$ & 27 & 4 \\
$(4,5)$ & 24 & $r=4$：$0+15+24=39$；$r=5$：$12+0+24=\mathbf{36}$ & 36 & 5 \\
$(1,3)$ & 27 & $r=1$：$0+28+27=55$；$r=2$：$9+8+27=\mathbf{44}$；$r=3$：$32+0+27=59$ & 44 & 2 \\
$(2,4)$ & 31 & $r=2$：$0+27+31=\mathbf{58}$；$r=3$：$16+12+31=59$；$r=4$：$28+0+31=59$ & 58 & 2 \\
$(3,5)$ & 31 & $r=3$：$0+36+31=67$；$r=4$：$8+15+31=\mathbf{54}$；$r=5$：$27+0+31=58$ & 54 & 4 \\
$(1,4)$ & 38 & $r=1$：$0+58+38=96$；$r=2$：$9+27+38=\mathbf{74}$；$r=3$：$32+12+38=82$；$r=4$：$44+0+38=82$ & 74 & 2 \\
$(2,5)$ & 43 & $r=2$：$0+54+43=97$；$r=3$：$16+36+43=95$；$r=4$：$28+15+43=\mathbf{86}$；$r=5$：$58+0+43=101$ & 86 & 4 \\
$(1,5)$ & 50 & $r=1$：$0+86+50=136$；$r=2$：$9+54+50=113$；$r=3$：$32+36+50=118$；$r=4$：$44+15+50=\mathbf{109}$；$r=5$：$74+0+50=124$ & 109 & 4 \\
\bottomrule
\end{tabular}
\end{center}
}

\begin{center}
\begin{minipage}{\linewidth}
\centering
\begin{tikzpicture}[x=0.85cm, y=0.62cm, baseline=(current bounding box.north)]
\node[font=\small] at (3.50,1.15) {$e[i,j]$};
\node[dp-hd] at (1,0.35) {$j=0$};
\node[dp-hd] at (2,0.35) {$1$};
\node[dp-hd] at (3,0.35) {$2$};
\node[dp-hd] at (4,0.35) {$3$};
\node[dp-hd] at (5,0.35) {$4$};
\node[dp-hd] at (6,0.35) {$5$};
\node[dp-hd, anchor=east] at (0.45,-0.5) {$i=1$};
\node[dp-c0, minimum width=0.85cm] at (1,-0.5) {0};
\node[dp-cd0, minimum width=0.85cm] at (2,-0.5) {9};
\node[dp-cd1, minimum width=0.85cm] at (3,-0.5) {32};
\node[dp-cd0, minimum width=0.85cm] at (4,-0.5) {44};
\node[dp-cd1, minimum width=0.85cm] at (5,-0.5) {74};
\node[dp-cd0, minimum width=0.85cm] at (6,-0.5) {109};
\node[dp-hd, anchor=east] at (0.45,-1.5) {$2$};
\node[dp-c0, minimum width=0.85cm] at (2,-1.5) {0};
\node[dp-cd0, minimum width=0.85cm] at (3,-1.5) {16};
\node[dp-cd1, minimum width=0.85cm] at (4,-1.5) {28};
\node[dp-cd0, minimum width=0.85cm] at (5,-1.5) {58};
\node[dp-cd1, minimum width=0.85cm] at (6,-1.5) {86};
\node[dp-hd, anchor=east] at (0.45,-2.5) {$3$};
\node[dp-c0, minimum width=0.85cm] at (3,-2.5) {0};
\node[dp-cd0, minimum width=0.85cm] at (4,-2.5) {8};
\node[dp-cd1, minimum width=0.85cm] at (5,-2.5) {27};
\node[dp-cd0, minimum width=0.85cm] at (6,-2.5) {54};
\node[dp-hd, anchor=east] at (0.45,-3.5) {$4$};
\node[dp-c0, minimum width=0.85cm] at (4,-3.5) {0};
\node[dp-cd0, minimum width=0.85cm] at (5,-3.5) {12};
\node[dp-cd1, minimum width=0.85cm] at (6,-3.5) {36};
\node[dp-hd, anchor=east] at (0.45,-4.5) {$5$};
\node[dp-c0, minimum width=0.85cm] at (5,-4.5) {0};
\node[dp-cd0, minimum width=0.85cm] at (6,-4.5) {15};
\node[dp-hd, anchor=east] at (0.45,-5.5) {$6$};
\node[dp-c0, minimum width=0.85cm] at (6,-5.5) {0};
\end{tikzpicture}\hspace{0.8cm}\begin{tikzpicture}[x=0.8cm, y=0.62cm, baseline=(current bounding box.north)]
\node[font=\small] at (3.00,1.15) {$\mathit{root}[i,j]$};
\node[dp-hd] at (1,0.35) {$j=1$};
\node[dp-hd] at (2,0.35) {$2$};
\node[dp-hd] at (3,0.35) {$3$};
\node[dp-hd] at (4,0.35) {$4$};
\node[dp-hd] at (5,0.35) {$5$};
\node[dp-hd, anchor=east] at (0.45,-0.5) {$i=1$};
\node[dp-cd0, minimum width=0.8cm] at (1,-0.5) {1};
\node[dp-cd1, minimum width=0.8cm] at (2,-0.5) {2};
\node[dp-cd0, minimum width=0.8cm] at (3,-0.5) {2};
\node[dp-cd1, minimum width=0.8cm] at (4,-0.5) {2};
\node[dp-cd0, minimum width=0.8cm] at (5,-0.5) {4};
\node[dp-hd, anchor=east] at (0.45,-1.5) {$2$};
\node[dp-cd0, minimum width=0.8cm] at (2,-1.5) {2};
\node[dp-cd1, minimum width=0.8cm] at (3,-1.5) {2};
\node[dp-cd0, minimum width=0.8cm] at (4,-1.5) {2};
\node[dp-cd1, minimum width=0.8cm] at (5,-1.5) {4};
\node[dp-hd, anchor=east] at (0.45,-2.5) {$3$};
\node[dp-cd0, minimum width=0.8cm] at (3,-2.5) {3};
\node[dp-cd1, minimum width=0.8cm] at (4,-2.5) {4};
\node[dp-cd0, minimum width=0.8cm] at (5,-2.5) {4};
\node[dp-hd, anchor=east] at (0.45,-3.5) {$4$};
\node[dp-cd0, minimum width=0.8cm] at (4,-3.5) {4};
\node[dp-cd1, minimum width=0.8cm] at (5,-3.5) {5};
\node[dp-hd, anchor=east] at (0.45,-4.5) {$5$};
\node[dp-cd0, minimum width=0.8cm] at (5,-4.5) {5};
\end{tikzpicture}
\diagcaption{$e$ 和 $\mathit{root}$ 兩張表（$w$ 的值見上表第二欄）。 灰底的是 $j = i - 1$ 的空樹，底色深淺交錯的是不
同長度的區間。 答案是 $e[1, 5] = 109$。}
\label{fig:dp-obst-tables}
\end{minipage}
\end{center}

\textbf{由 root 表建樹}：$\mathit{root}[1, 5] = 4$，根是 $k_4 = \texttt{int}$；左邊是
$k_1..k_3$，$\mathit{root}[1, 3] = 2$，根是 \texttt{for}，它的左右各是 $\mathit{root}[1,1] = 1$
的 \texttt{char} 和 $\mathit{root}[3,3] = 3$ 的 \texttt{goto}；右邊是 $k_5..k_5$，就是
\texttt{while}。 這正是圖~\ref{fig:dp-obst-trees}(b)。 和矩陣鏈乘的 \proc{Print-Parens} 一
樣，是一個 $O(n)$ 的遞迴。

表裡有一個現象：每一列的 $\mathit{root}$ 由左往右不會變小，每一欄由上往下也不會變小。 Knuth
在 1971 年證明這對任何次數都成立：
\[
\mathit{root}[i, j - 1] \le \mathit{root}[i, j] \le \mathit{root}[i + 1, j] ,
\]
所以第 10 行的 $r$ 只要從 $\mathit{root}[i, j-1]$ 試到 $\mathit{root}[i+1, j]$，共
$\mathit{root}[i+1, j] - \mathit{root}[i, j-1] + 1$ 個。 固定區間長度 $\ell$，把 $i = 1, 2,
\ldots, n - \ell + 1$ 這些個數加起來（$j = i + \ell - 1$ 跟著 $i$ 變）：第 $i$ 項加的
$\mathit{root}[i+1, j]$，正好是第 $i + 1$ 項減的 $\mathit{root}[(i+1), (j+1) - 1]$，兩兩抵銷，最
後只剩 $\mathit{root}[n - \ell + 2, n] - \mathit{root}[1, \ell - 1]$ 和 $n - \ell + 1$ 個 $1$，不超過
$n + (n - \ell + 1) \le 2n$。 每個 $\ell$ 是 $O(n)$，$n$ 個 $\ell$ 合起來，總時間
降到 $\Theta(n^2)$。 例如 $e[1, 5]$ 只要試
$r = \mathit{root}[1, 4] = 2$ 到 $\mathit{root}[2, 5] = 4$ 這三個。

\begin{examnote}
最佳二元搜尋樹的題目，先看清楚題目的成本怎麼算。 本書數「和鍵比較的次數」：鍵的成本是深度
加一，外部節點的成本是深度，所以空樹 $e[i, i-1] = 0$。 CLRS 把走到外部節點也算一次，外部
節點的成本是深度加一，所以它的 $e[i, i-1] = q_{i-1}$，每一格的 $e$ 都比本書多
$q_{i-1} + q_i + \cdots + q_j$；本節的例子照 CLRS 的成本模型是 $109 + 19 = 128$。 另一種常見
的寫法（例如 Horowitz、Sahni 等人的 \textit{Computer Algorithms}）用 $c(i, j)$ 表示
$a_{i+1}, \ldots, a_j$ 這一段，$c(i, i) = 0$、$w(i, i) = q_i$，數法和本書相同，只是索引差一。 幾種成本模型的 $\mathit{root}$ 表和最佳樹都一樣（相差的是一個和
樹的形狀無關的常數），但 $e$ 的數值不同，答案要照題目的定義寫，必要時寫一句「依題意，查
不到時不計走到外部節點那一步」。
\end{examnote}

\section{0/1 背包}
\label{sec:dp-knapsack}

\subsection{問題}

期末專題還剩 $10$ 天。 組員列出五個還沒做的功能，每個功能要花的天數 $w_i$ 和做完能加的分數
$v_i$ 如下；功能只能整個做完或整個不做，總天數不能超過 $10$ 天。 要做哪幾個，分數最高？

\begin{center}
\begin{tabular}{c*{5}{c}}
\toprule
功能 $i$ & 1 & 2 & 3 & 4 & 5 \\
\midrule
天數 $w_i$ & 3 & 2 & 4 & 5 & 2 \\
分數 $v_i$ & 10 & 7 & 11 & 13 & 6 \\
每天的分數 $v_i / w_i$ & 3.33 & 3.5 & 2.75 & 2.6 & 3 \\
\bottomrule
\end{tabular}
\end{center}

一般地，\textbf{0/1 背包問題}（0/1 knapsack problem）：有 $n$ 件物品，第 $i$ 件重 $w_i$、價值
$v_i$（都是正整數），背包容量 $W$。 每件物品拿或不拿（所以叫 0/1），總重量不超過 $W$，求
總價值最大。

先試一個直覺的做法：照「每天的分數」由高到低挑，放得下就做。 依序是功能 2、1、5（共 $7$
天、$23$ 分），剩下 $3$ 天，功能 3 和 4 都放不下，得到 $23$ 分。 可是功能 1、2、4 剛好
$10$ 天、$30$ 分。 照比值挑的貪心法對 0/1 背包不對；物品可以只拿一部分時（分數背包），它
才是對的。 第~19~章會證明分數背包的貪心法，並把兩者放在一起比較。 全部試一遍要 $2^n$ 種組
合，下面用動態規劃。

\subsection{遞迴式}

看最後一件物品 $n$。 最佳解不是不拿它，就是拿它：
\begin{itemize}
\item 不拿：剩下的就是「前 $n - 1$ 件、容量 $W$」的最佳解。
\item 拿（需要 $w_n \le W$）：其餘的物品要在容量 $W - w_n$ 裡挑，而且一定是「前 $n - 1$ 件、
容量 $W - w_n$」的最佳解，否則換成更好的那組，整體就更好（剪下貼上）。
\end{itemize}
子問題要同時記住「考慮到第幾件」和「剩多少容量」。 令 $K[i, c]$ 為只用前 $i$ 件物品、容量
$c$ 時的最大價值，$0 \le i \le n$，$0 \le c \le W$：
\[
K[i, c] = \begin{cases}
0, & i = 0, \\
K[i - 1, c], & i \ge 1 \text{ 且 } w_i > c, \\
\max\bigl(K[i - 1, c],\ K[i - 1, c - w_i] + v_i\bigr), & i \ge 1 \text{ 且 } w_i \le c .
\end{cases}
\]
答案是 $K[n, W]$。 第 $i$ 列只用到第 $i - 1$ 列，所以一列一列往下填。

\begin{pseudo}[label=alg:dp-knapsack]{0/1 背包的填表}
\Function{Knapsack}{$w, v, n, W$}
    \For{$c \gets 0$ \To $W$}
        \State $K[0, c] \gets 0$
    \EndFor
    \For{$i \gets 1$ \To $n$}
        \For{$c \gets 0$ \To $W$}
            \State $K[i, c] \gets K[i - 1, c]$ \Comment{不拿第 $i$ 件}
            \If{$w_i \le c$ \AND $K[i - 1, c - w_i] + v_i > K[i, c]$}
                \State $K[i, c] \gets K[i - 1, c - w_i] + v_i$ \Comment{拿第 $i$ 件}
            \EndIf
        \EndFor
    \EndFor
    \State \Return $K$
\EndFunction
\end{pseudo}

\textbf{還原}：從 $K[n, W]$ 往上看。 $K[i, c] \ne K[i - 1, c]$ 表示第 $i$ 件一定有拿（不拿的
話值會和上一列一樣），於是容量減掉 $w_i$ 再看上一列；相等就表示不拿第 $i$ 件也能達到，照第 7
行的 $>$，本書這時當作不拿。

\begin{pseudo}[label=alg:dp-knapsack-items]{還原要拿的物品}
\Procedure{Knapsack-Items}{$K, w, n, W$}
    \State $c \gets W$
    \For{$i \gets n$ \DownTo $1$}
        \If{$K[i, c] \ne K[i - 1, c]$}
            \State 印出 $i$ \Comment{第 $i$ 件有拿}
            \State $c \gets c - w_i$
        \EndIf
    \EndFor
\EndProcedure
\end{pseudo}

\subsection{手算追蹤}
\label{sec:dp-knapsack-trace}

\begin{center}
\begin{minipage}{\linewidth}
\centering
\begin{tikzpicture}[x=0.72cm, y=0.6cm]
\node[dp-hd, anchor=east] at (-0.5,0.9) {$c$};
\node[dp-hd] at (0,0.9) {0};
\node[dp-hd] at (1,0.9) {1};
\node[dp-hd] at (2,0.9) {2};
\node[dp-hd] at (3,0.9) {3};
\node[dp-hd] at (4,0.9) {4};
\node[dp-hd] at (5,0.9) {5};
\node[dp-hd] at (6,0.9) {6};
\node[dp-hd] at (7,0.9) {7};
\node[dp-hd] at (8,0.9) {8};
\node[dp-hd] at (9,0.9) {9};
\node[dp-hd] at (10,0.9) {10};
\node[dp-hd, anchor=east] at (-0.45,0) {$i=0$};
\node[dp-cpath, minimum width=0.72cm, minimum height=0.6cm] at (0,0) {0};
\node[dp-c, minimum width=0.72cm, minimum height=0.6cm] at (1,0) {0};
\node[dp-c, minimum width=0.72cm, minimum height=0.6cm] at (2,0) {0};
\node[dp-c, minimum width=0.72cm, minimum height=0.6cm] at (3,0) {0};
\node[dp-c, minimum width=0.72cm, minimum height=0.6cm] at (4,0) {0};
\node[dp-c, minimum width=0.72cm, minimum height=0.6cm] at (5,0) {0};
\node[dp-c, minimum width=0.72cm, minimum height=0.6cm] at (6,0) {0};
\node[dp-c, minimum width=0.72cm, minimum height=0.6cm] at (7,0) {0};
\node[dp-c, minimum width=0.72cm, minimum height=0.6cm] at (8,0) {0};
\node[dp-c, minimum width=0.72cm, minimum height=0.6cm] at (9,0) {0};
\node[dp-c, minimum width=0.72cm, minimum height=0.6cm] at (10,0) {0};
\node[dp-hd, anchor=east] at (-0.45,-1) {$i=1$ (3,10)};
\node[dp-c, minimum width=0.72cm, minimum height=0.6cm] at (0,-1) {0};
\node[dp-c, minimum width=0.72cm, minimum height=0.6cm] at (1,-1) {0};
\node[dp-c, minimum width=0.72cm, minimum height=0.6cm] at (2,-1) {0};
\node[dp-chit, minimum width=0.72cm, minimum height=0.6cm] at (3,-1) {10};
\node[dp-c, minimum width=0.72cm, minimum height=0.6cm] at (4,-1) {10};
\node[dp-c, minimum width=0.72cm, minimum height=0.6cm] at (5,-1) {10};
\node[dp-c, minimum width=0.72cm, minimum height=0.6cm] at (6,-1) {10};
\node[dp-c, minimum width=0.72cm, minimum height=0.6cm] at (7,-1) {10};
\node[dp-c, minimum width=0.72cm, minimum height=0.6cm] at (8,-1) {10};
\node[dp-c, minimum width=0.72cm, minimum height=0.6cm] at (9,-1) {10};
\node[dp-c, minimum width=0.72cm, minimum height=0.6cm] at (10,-1) {10};
\node[dp-hd, anchor=east] at (-0.45,-2) {$i=2$ (2,7)};
\node[dp-c, minimum width=0.72cm, minimum height=0.6cm] at (0,-2) {0};
\node[dp-c, minimum width=0.72cm, minimum height=0.6cm] at (1,-2) {0};
\node[dp-c, minimum width=0.72cm, minimum height=0.6cm] at (2,-2) {7};
\node[dp-c, minimum width=0.72cm, minimum height=0.6cm] at (3,-2) {10};
\node[dp-c, minimum width=0.72cm, minimum height=0.6cm] at (4,-2) {10};
\node[dp-chit, minimum width=0.72cm, minimum height=0.6cm] at (5,-2) {17};
\node[dp-c, minimum width=0.72cm, minimum height=0.6cm] at (6,-2) {17};
\node[dp-c, minimum width=0.72cm, minimum height=0.6cm] at (7,-2) {17};
\node[dp-c, minimum width=0.72cm, minimum height=0.6cm] at (8,-2) {17};
\node[dp-c, minimum width=0.72cm, minimum height=0.6cm] at (9,-2) {17};
\node[dp-c, minimum width=0.72cm, minimum height=0.6cm] at (10,-2) {17};
\node[dp-hd, anchor=east] at (-0.45,-3) {$i=3$ (4,11)};
\node[dp-c, minimum width=0.72cm, minimum height=0.6cm] at (0,-3) {0};
\node[dp-c, minimum width=0.72cm, minimum height=0.6cm] at (1,-3) {0};
\node[dp-c, minimum width=0.72cm, minimum height=0.6cm] at (2,-3) {7};
\node[dp-c, minimum width=0.72cm, minimum height=0.6cm] at (3,-3) {10};
\node[dp-c, minimum width=0.72cm, minimum height=0.6cm] at (4,-3) {11};
\node[dp-cpath, minimum width=0.72cm, minimum height=0.6cm] at (5,-3) {17};
\node[dp-c, minimum width=0.72cm, minimum height=0.6cm] at (6,-3) {18};
\node[dp-c, minimum width=0.72cm, minimum height=0.6cm] at (7,-3) {21};
\node[dp-c, minimum width=0.72cm, minimum height=0.6cm] at (8,-3) {21};
\node[dp-c, minimum width=0.72cm, minimum height=0.6cm] at (9,-3) {28};
\node[dp-c, minimum width=0.72cm, minimum height=0.6cm] at (10,-3) {28};
\node[dp-hd, anchor=east] at (-0.45,-4) {$i=4$ (5,13)};
\node[dp-c, minimum width=0.72cm, minimum height=0.6cm] at (0,-4) {0};
\node[dp-c, minimum width=0.72cm, minimum height=0.6cm] at (1,-4) {0};
\node[dp-c, minimum width=0.72cm, minimum height=0.6cm] at (2,-4) {7};
\node[dp-c, minimum width=0.72cm, minimum height=0.6cm] at (3,-4) {10};
\node[dp-c, minimum width=0.72cm, minimum height=0.6cm] at (4,-4) {11};
\node[dp-c, minimum width=0.72cm, minimum height=0.6cm] at (5,-4) {17};
\node[dp-c, minimum width=0.72cm, minimum height=0.6cm] at (6,-4) {18};
\node[dp-c, minimum width=0.72cm, minimum height=0.6cm] at (7,-4) {21};
\node[dp-c, minimum width=0.72cm, minimum height=0.6cm] at (8,-4) {23};
\node[dp-c, minimum width=0.72cm, minimum height=0.6cm] at (9,-4) {28};
\node[dp-chit, minimum width=0.72cm, minimum height=0.6cm] at (10,-4) {30};
\node[dp-hd, anchor=east] at (-0.45,-5) {$i=5$ (2,6)};
\node[dp-c, minimum width=0.72cm, minimum height=0.6cm] at (0,-5) {0};
\node[dp-c, minimum width=0.72cm, minimum height=0.6cm] at (1,-5) {0};
\node[dp-c, minimum width=0.72cm, minimum height=0.6cm] at (2,-5) {7};
\node[dp-c, minimum width=0.72cm, minimum height=0.6cm] at (3,-5) {10};
\node[dp-c, minimum width=0.72cm, minimum height=0.6cm] at (4,-5) {13};
\node[dp-c, minimum width=0.72cm, minimum height=0.6cm] at (5,-5) {17};
\node[dp-c, minimum width=0.72cm, minimum height=0.6cm] at (6,-5) {18};
\node[dp-c, minimum width=0.72cm, minimum height=0.6cm] at (7,-5) {23};
\node[dp-c, minimum width=0.72cm, minimum height=0.6cm] at (8,-5) {24};
\node[dp-c, minimum width=0.72cm, minimum height=0.6cm] at (9,-5) {28};
\node[dp-cpath, minimum width=0.72cm, minimum height=0.6cm] at (10,-5) {30};
\end{tikzpicture}
\diagcaption{專題例子的 $K$ 表。 列左邊的括號是 $(w_i, v_i)$。 從 $K[5, 10] = 30$ 往上還原，
經過的格子加框；粉紅色的格子和上一列不同，表示拿了那一列的物品：功能 4、2、1。}
\label{fig:dp-knapsack-table}
\end{minipage}
\end{center}

幾格是這樣算的：
\begin{itemize}
\item $K[2, 5]$：功能 2 重 $2$，$\max(K[1, 5], K[1, 3] + 7) = \max(10, 17) = 17$，也就是功能
1、2 都做。
\item $K[3, 7]$：$\max(K[2, 7], K[2, 3] + 11) = \max(17, 21) = 21$，做功能 1、3。
\item $K[4, 10]$：$\max(K[3, 10], K[3, 5] + 13) = \max(28, 30) = 30$。
\item $K[5, 10]$：$\max(K[4, 10], K[4, 8] + 6) = \max(30, 29) = 30$，不做功能 5。
\end{itemize}

還原：$K[5, 10] = K[4, 10]$，不做功能 5；$K[4, 10] = 30 \ne K[3, 10] = 28$，做功能 4，容量剩
$5$；$K[3, 5] = K[2, 5] = 17$，不做功能 3；$K[2, 5] = 17 \ne K[1, 5] = 10$，做功能 2，剩 $3$；
$K[1, 3] = 10 \ne K[0, 3] = 0$，做功能 1，剩 $0$。 答案是功能 1、2、4，$10$ 天、$30$ 分。

\subsection{一維滾動陣列}
\label{sec:dp-knapsack-1d}

第 $i$ 列只用到第 $i - 1$ 列，所以和第~\ref{sec:dp-bottomup}~節的巴斯卡三角形一樣，可以只留一
列 $d[0..W]$，就地從第 $i - 1$ 列改成第 $i$ 列。 更新 $d[c]$ 時要用的是\textbf{舊的}
$d[c - w_i]$（第 $i - 1$ 列的值），而 $c - w_i < c$，所以容量要\textbf{由大到小}跑：改 $d[c]$
的時候，比它小的格子都還沒動過。

\begin{pseudo}[label=alg:dp-knapsack-1d]{0/1 背包，一維陣列}
\Function{Knapsack-1D}{$w, v, n, W$}
    \For{$c \gets 0$ \To $W$}
        \State $d[c] \gets 0$
    \EndFor
    \For{$i \gets 1$ \To $n$}
        \For{$c \gets W$ \DownTo $w_i$}
            \State $d[c] \gets \max(d[c],\ d[c - w_i] + v_i)$
        \EndFor
    \EndFor
    \State \Return $d[W]$
\EndFunction
\end{pseudo}

$c < w_i$ 的格子放不下第 $i$ 件，值和上一列一樣，所以第 5 行的迴圈只跑到 $w_i$。 空間從
$\Theta(nW)$ 降到 $\Theta(W)$，時間不變；但和 LCS 只留兩列一樣，表被蓋掉之後就不能用
\proc{Knapsack-Items} 還原了。

如果第 5 行寫成 $c \gets w_i$ \To $W$ 由小到大，算 $d[c]$ 時 $d[c - w_i]$ 可能已經是這一列的
新值，也就是\textbf{已經拿過第 $i$ 件}的結果，於是同一件物品會被拿好幾次。 只看功能 2：
$d[2] = 7$，接著 $d[4] = d[2] + 7 = 14$、$d[6] = 21$、……、$d[10] = 35$，等於把功能 2 做了
五次。 這個「錯誤」的版本其實正好解另一個問題：每種物品可以拿\textbf{任意多個}的
\textbf{無限背包}（unbounded knapsack），它的遞迴式是
\[
U[c] = \max\Bigl(0,\ \max_{i : w_i \le c} \bigl(U[c - w_i] + v_i\bigr)\Bigr),
\]
由小到大算 $U[c]$ 時，$U[c - w_i]$ 本來就該是已經考慮過所有物品（包括第 $i$ 件）的值。 兩
個問題只差一個迴圈的方向，這一點很容易考，也很容易寫錯。

\subsection{擬多項式時間}
\label{sec:dp-pseudo-poly}

\proc{Knapsack} 的表有 $(n + 1)(W + 1)$ 格，每格 $O(1)$，時間 $\Theta(nW)$。 這看起來是多項式，
但第~\ref{sec:ana-input-size}~節說過，輸入的大小要用寫下輸入所需的位元數來量。 $W$
這個數只要 $b = \lceil \lg(W + 1) \rceil$ 個位元就寫得下，所以 $\Theta(nW) = \Theta(n 2^b)$
是輸入長度的\textbf{指數}函數。 容量從 $10$ 改成 $10^{9}$，輸入只多了幾個字元，表卻大了一億
倍。 這種「對輸入的數值是多項式、對位元數是指數」的時間，就是\textbf{擬多項式時間}。

$W$ 不大時（例如練習題的 $W \le 10^5$），$\Theta(nW)$ 非常好用；$W$ 很大時這張表就開不出來。
0/1 背包的決策版本是 NP 完備問題（第~28~章），目前沒有已知的多項式時間演算法，所以擬多項
式已經是很好的結果。 另一種 DP 改用「達到價值 $v$ 最少要多重」當子問題，時間是
$\Theta(n \sum v_i)$，價值小、重量大時比較好；第~29~章的近似演算法就從這個版本出發。 第~20~
章則用分支限界在不填表的情況下搜尋。

\subsection{C 程式}

\texttt{knap\_2d} 填完整的表並還原；\texttt{knap\_1d} 用一維陣列，參數
\texttt{down} 決定容量的方向。 為了和虛擬碼對照，物品從 $1$ 開始編號，\texttt{w[0]}、
\texttt{v[0]} 不用。

\begin{example}[label=ex:dp-knapsack]{c}{(knapsack.c) 0/1 背包的二維表與一維陣列}
#include <stdio.h>
#include <stdlib.h>

#define MAXN 100

// 二維表 K[i][c]：只用前 i 件、容量 c 的最大價值。take[i] 記錄有沒有拿第 i 件。
long long knap_2d(int n, int cap, const int *w, const int *v, int *take)
{
    long long *K = calloc((size_t)(n + 1) * (cap + 1), sizeof *K);
    if (K == NULL) {
        return -1;
    }
#define T(i, c) K[(i) * (cap + 1) + (c)]
    for (int i = 1; i <= n; i++) {
        for (int c = 0; c <= cap; c++) {
            T(i, c) = T(i - 1, c);                      // 不拿第 i 件
            if (w[i] <= c && T(i - 1, c - w[i]) + v[i] > T(i, c)) {
                T(i, c) = T(i - 1, c - w[i]) + v[i];    // 拿第 i 件
            }
        }
    }
    long long best = T(n, cap);
    for (int i = n, c = cap; i >= 1; i--) {             // 從最後一列往上還原
        take[i] = (T(i, c) != T(i - 1, c));
        if (take[i]) {
            c -= w[i];
        }
    }
#undef T
    free(K);
    return best;
}

// 一維陣列 d[c]。down = 1 時容量由大到小（0/1 背包），0 時由小到大。
long long knap_1d(int n, int cap, const int *w, const int *v, int down)
{
    long long *d = calloc(cap + 1, sizeof *d);
    if (d == NULL) {
        return -1;
    }
    for (int i = 1; i <= n; i++) {
        for (int t = 0; t <= cap - w[i]; t++) {
            int c = down ? cap - t : w[i] + t;
            if (d[c - w[i]] + v[i] > d[c]) {
                d[c] = d[c - w[i]] + v[i];
            }
        }
    }
    long long best = d[cap];
    free(d);
    return best;
}

int main(void)
{
    int n, cap, w[MAXN + 1], v[MAXN + 1], take[MAXN + 1];
    if (scanf("%d %d", &n, &cap) != 2 || n < 1 || n > MAXN || cap < 0) {
        return 1;
    }
    for (int i = 1; i <= n; i++) {
        if (scanf("%d", &w[i]) != 1) {
            return 1;
        }
    }
    for (int i = 1; i <= n; i++) {
        if (scanf("%d", &v[i]) != 1) {
            return 1;
        }
    }
    printf("2D table: best = %lld, items =", knap_2d(n, cap, w, v, take));
    for (int i = 1; i <= n; i++) {
        if (take[i]) {
            printf(" %d", i);
        }
    }
    printf("\n1D, capacity downward: best = %lld\n", knap_1d(n, cap, w, v, 1));
    printf("1D, capacity upward: best = %lld\n", knap_1d(n, cap, w, v, 0));
    return 0;
}
\end{example}
\iosample{5 10 \\ 3 2 4 5 2 \\ 10 7 11 13 6}{2D table: best = 30, items = 1 2 4 \\ 1D, capacity downward: best = 30 \\ 1D, capacity upward: best = 35}

由大到小的一維版本和二維表一樣是 $30$；由小到大的版本得到 $35$，是把功能 2 做五次的無限背
包答案，對 0/1 背包來說是錯的。 一維迴圈為了同時支援兩個方向，用 \texttt{t} 從 $0$ 數到
$W - w_i$，再換算成容量 \texttt{c}；只寫 0/1 背包的話，直接寫
\texttt{for (int c = cap; c >= w[i]; c--)} 就好。

\section{編輯距離與最長遞增子序列}
\label{sec:dp-edit-lis}

這一節的兩個問題，子問題都是「前綴」，和 LCS 是同一類。

\subsection{編輯距離}
\label{sec:dp-edit}

拼字檢查要找「和打錯的字最接近的正確單字」，就要量兩個字串有多接近。 把字串 $X$ 改成 $Y$，
每一步可以做三種操作之一：
\begin{itemize}
\item \textbf{插入}一個字元；
\item \textbf{刪除}一個字元；
\item 把一個字元\textbf{替換}成另一個字元。
\end{itemize}
每種操作的成本都是 $1$，最少的操作次數叫 $X$ 到 $Y$ 的\textbf{編輯距離}（edit distance），也
叫 Levenshtein 距離。 例如 \texttt{PARTNER} 改成 \texttt{PRINTER}：刪掉 \texttt{A}、把
\texttt{T} 換成 \texttt{I}、在 \texttt{N} 後面插入 \texttt{T}，三步。 不可能更少，下面的表會
證明這件事。

把改的過程寫成上下兩列的\textbf{對齊}比較好想：
\begin{center}
\texttt{P A R T N - E R}\\
\texttt{P - R I N T E R}
\end{center}
每一行是三種之一：上下都有字元（相同就不花成本，不同是一次替換）、只有上面（刪除）、只有下面
（插入）。 編輯距離就是最好的對齊裡「不是相同字元」的行數。 看對齊的\textbf{最後一行}：
\begin{itemize}
\item 上面是 $x_m$、下面是 $y_n$：前面是 $X_{m-1}$ 和 $Y_{n-1}$ 的對齊，成本加 $0$ 或 $1$；
\item 上面是 $x_m$、下面空著（刪除 $x_m$）：前面是 $X_{m-1}$ 和 $Y_n$ 的對齊，成本加 $1$；
\item 上面空著、下面是 $y_n$（插入 $y_n$）：前面是 $X_m$ 和 $Y_{n-1}$ 的對齊，成本加 $1$。
\end{itemize}
前面那一段當然要用最好的對齊（剪下貼上）。 令 $D[i, j]$ 為 $X_i$ 改成 $Y_j$ 的編輯距離：
\[
D[i, j] = \min\begin{cases}
D[i - 1, j - 1] + \delta(x_i, y_j), \\
D[i - 1, j] + 1, \\
D[i, j - 1] + 1,
\end{cases}
\qquad
\delta(x_i, y_j) = \begin{cases} 0, & x_i = y_j, \\ 1, & x_i \ne y_j, \end{cases}
\]
邊界是 $D[i, 0] = i$（刪掉 $i$ 個字元）、$D[0, j] = j$（插入 $j$ 個字元）。 和 LCS 不同，
$x_i = y_j$ 時三項還是要一起比；不過可以證明這時取左上角一定不會比較差（觀念題
第~\ref{dp-q-edit-match}~題）。

\begin{pseudo}[label=alg:dp-edit]{編輯距離}
\Function{Edit-Distance}{$X, Y, m, n$}
    \For{$i \gets 0$ \To $m$}
        \State $D[i, 0] \gets i$
    \EndFor
    \For{$j \gets 0$ \To $n$}
        \State $D[0, j] \gets j$
    \EndFor
    \For{$i \gets 1$ \To $m$}
        \For{$j \gets 1$ \To $n$}
            \State $\delta \gets 0$；若 $x_i \ne y_j$，$\delta \gets 1$
            \State $D[i, j] \gets \min(D[i - 1, j - 1] + \delta,\ D[i - 1, j] + 1,\ D[i, j - 1] + 1)$
        \EndFor
    \EndFor
    \State \Return $D[m, n]$
\EndFunction
\end{pseudo}

時間和空間都是 $\Theta(mn)$；只要距離的話，同樣可以只留兩列。

\begin{center}
\begin{minipage}{\linewidth}
\centering
\begin{tikzpicture}[x=0.78cm, y=0.6cm]
\node[dp-hd] at (0,1.55) {0};
\node[dp-hd] at (1,1.55) {1};
\node[dp-hd] at (1,0.85) {\texttt{P}};
\node[dp-hd] at (2,1.55) {2};
\node[dp-hd] at (2,0.85) {\texttt{R}};
\node[dp-hd] at (3,1.55) {3};
\node[dp-hd] at (3,0.85) {\texttt{I}};
\node[dp-hd] at (4,1.55) {4};
\node[dp-hd] at (4,0.85) {\texttt{N}};
\node[dp-hd] at (5,1.55) {5};
\node[dp-hd] at (5,0.85) {\texttt{T}};
\node[dp-hd] at (6,1.55) {6};
\node[dp-hd] at (6,0.85) {\texttt{E}};
\node[dp-hd] at (7,1.55) {7};
\node[dp-hd] at (7,0.85) {\texttt{R}};
\node[dp-hd] at (-1.2,0) {0};
\node[dp-cpath, minimum width=0.78cm, minimum height=0.6cm] at (0,0) {0};
\node[dp-c0, minimum width=0.78cm, minimum height=0.6cm] at (1,0) {1};
\node[dp-c0, minimum width=0.78cm, minimum height=0.6cm] at (2,0) {2};
\node[dp-c0, minimum width=0.78cm, minimum height=0.6cm] at (3,0) {3};
\node[dp-c0, minimum width=0.78cm, minimum height=0.6cm] at (4,0) {4};
\node[dp-c0, minimum width=0.78cm, minimum height=0.6cm] at (5,0) {5};
\node[dp-c0, minimum width=0.78cm, minimum height=0.6cm] at (6,0) {6};
\node[dp-c0, minimum width=0.78cm, minimum height=0.6cm] at (7,0) {7};
\node[dp-hd] at (-1.2,-1) {1};
\node[dp-hd] at (-0.65,-1) {\texttt{P}};
\node[dp-c0, minimum width=0.78cm, minimum height=0.6cm] at (0,-1) {1};
\node[dp-cpath, minimum width=0.78cm, minimum height=0.6cm] at (1,-1) {0};
\node[dp-c, minimum width=0.78cm, minimum height=0.6cm] at (2,-1) {1};
\node[dp-c, minimum width=0.78cm, minimum height=0.6cm] at (3,-1) {2};
\node[dp-c, minimum width=0.78cm, minimum height=0.6cm] at (4,-1) {3};
\node[dp-c, minimum width=0.78cm, minimum height=0.6cm] at (5,-1) {4};
\node[dp-c, minimum width=0.78cm, minimum height=0.6cm] at (6,-1) {5};
\node[dp-c, minimum width=0.78cm, minimum height=0.6cm] at (7,-1) {6};
\node[dp-hd] at (-1.2,-2) {2};
\node[dp-hd] at (-0.65,-2) {\texttt{A}};
\node[dp-c0, minimum width=0.78cm, minimum height=0.6cm] at (0,-2) {2};
\node[dp-cpath, minimum width=0.78cm, minimum height=0.6cm] at (1,-2) {1};
\node[dp-c, minimum width=0.78cm, minimum height=0.6cm] at (2,-2) {1};
\node[dp-c, minimum width=0.78cm, minimum height=0.6cm] at (3,-2) {2};
\node[dp-c, minimum width=0.78cm, minimum height=0.6cm] at (4,-2) {3};
\node[dp-c, minimum width=0.78cm, minimum height=0.6cm] at (5,-2) {4};
\node[dp-c, minimum width=0.78cm, minimum height=0.6cm] at (6,-2) {5};
\node[dp-c, minimum width=0.78cm, minimum height=0.6cm] at (7,-2) {6};
\node[dp-hd] at (-1.2,-3) {3};
\node[dp-hd] at (-0.65,-3) {\texttt{R}};
\node[dp-c0, minimum width=0.78cm, minimum height=0.6cm] at (0,-3) {3};
\node[dp-c, minimum width=0.78cm, minimum height=0.6cm] at (1,-3) {2};
\node[dp-cpath, minimum width=0.78cm, minimum height=0.6cm] at (2,-3) {1};
\node[dp-c, minimum width=0.78cm, minimum height=0.6cm] at (3,-3) {2};
\node[dp-c, minimum width=0.78cm, minimum height=0.6cm] at (4,-3) {3};
\node[dp-c, minimum width=0.78cm, minimum height=0.6cm] at (5,-3) {4};
\node[dp-c, minimum width=0.78cm, minimum height=0.6cm] at (6,-3) {5};
\node[dp-c, minimum width=0.78cm, minimum height=0.6cm] at (7,-3) {5};
\node[dp-hd] at (-1.2,-4) {4};
\node[dp-hd] at (-0.65,-4) {\texttt{T}};
\node[dp-c0, minimum width=0.78cm, minimum height=0.6cm] at (0,-4) {4};
\node[dp-c, minimum width=0.78cm, minimum height=0.6cm] at (1,-4) {3};
\node[dp-c, minimum width=0.78cm, minimum height=0.6cm] at (2,-4) {2};
\node[dp-cpath, minimum width=0.78cm, minimum height=0.6cm] at (3,-4) {2};
\node[dp-c, minimum width=0.78cm, minimum height=0.6cm] at (4,-4) {3};
\node[dp-c, minimum width=0.78cm, minimum height=0.6cm] at (5,-4) {3};
\node[dp-c, minimum width=0.78cm, minimum height=0.6cm] at (6,-4) {4};
\node[dp-c, minimum width=0.78cm, minimum height=0.6cm] at (7,-4) {5};
\node[dp-hd] at (-1.2,-5) {5};
\node[dp-hd] at (-0.65,-5) {\texttt{N}};
\node[dp-c0, minimum width=0.78cm, minimum height=0.6cm] at (0,-5) {5};
\node[dp-c, minimum width=0.78cm, minimum height=0.6cm] at (1,-5) {4};
\node[dp-c, minimum width=0.78cm, minimum height=0.6cm] at (2,-5) {3};
\node[dp-c, minimum width=0.78cm, minimum height=0.6cm] at (3,-5) {3};
\node[dp-cpath, minimum width=0.78cm, minimum height=0.6cm] at (4,-5) {2};
\node[dp-cpath, minimum width=0.78cm, minimum height=0.6cm] at (5,-5) {3};
\node[dp-c, minimum width=0.78cm, minimum height=0.6cm] at (6,-5) {4};
\node[dp-c, minimum width=0.78cm, minimum height=0.6cm] at (7,-5) {5};
\node[dp-hd] at (-1.2,-6) {6};
\node[dp-hd] at (-0.65,-6) {\texttt{E}};
\node[dp-c0, minimum width=0.78cm, minimum height=0.6cm] at (0,-6) {6};
\node[dp-c, minimum width=0.78cm, minimum height=0.6cm] at (1,-6) {5};
\node[dp-c, minimum width=0.78cm, minimum height=0.6cm] at (2,-6) {4};
\node[dp-c, minimum width=0.78cm, minimum height=0.6cm] at (3,-6) {4};
\node[dp-c, minimum width=0.78cm, minimum height=0.6cm] at (4,-6) {3};
\node[dp-c, minimum width=0.78cm, minimum height=0.6cm] at (5,-6) {3};
\node[dp-cpath, minimum width=0.78cm, minimum height=0.6cm] at (6,-6) {3};
\node[dp-c, minimum width=0.78cm, minimum height=0.6cm] at (7,-6) {4};
\node[dp-hd] at (-1.2,-7) {7};
\node[dp-hd] at (-0.65,-7) {\texttt{R}};
\node[dp-c0, minimum width=0.78cm, minimum height=0.6cm] at (0,-7) {7};
\node[dp-c, minimum width=0.78cm, minimum height=0.6cm] at (1,-7) {6};
\node[dp-c, minimum width=0.78cm, minimum height=0.6cm] at (2,-7) {5};
\node[dp-c, minimum width=0.78cm, minimum height=0.6cm] at (3,-7) {5};
\node[dp-c, minimum width=0.78cm, minimum height=0.6cm] at (4,-7) {4};
\node[dp-c, minimum width=0.78cm, minimum height=0.6cm] at (5,-7) {4};
\node[dp-c, minimum width=0.78cm, minimum height=0.6cm] at (6,-7) {4};
\node[dp-cpath, minimum width=0.78cm, minimum height=0.6cm] at (7,-7) {3};
\end{tikzpicture}
\diagcaption{\texttt{PARTNER}（列）改成 \texttt{PRINTER}（欄）的 $D$ 表，答案 $D[7, 7] = 3$。
加框的是還原時走過的格子。}
\label{fig:dp-edit-table}
\end{minipage}
\end{center}

\textbf{還原。} 從 $[m, n]$ 往回走，每一步看這一格的值是由哪一項得到的，依序檢查：左上（相同
或替換）、上面（刪除 $x_i$）、左邊（插入 $y_j$），第一個符合的就走。 這組字串：
$[7,7] \to [6,6] \to [5,5]$ 是兩個相同字元 \texttt{R}、\texttt{E}；$[5,5]$ 的值 $3$ 不等於
$D[4,4] + 1 = 4$（\texttt{N} 和 \texttt{T} 不同），也不等於 $D[4,5] + 1 = 4$，等於
$D[5,4] + 1 = 3$，所以是插入 $y_5 = \texttt{T}$，走到 $[5, 4]$；接著 \texttt{N} 相同到 $[4,3]$；
$[4,3]$ 是 \texttt{T} 換成 \texttt{I}（$D[3,2] + 1 = 2$）；\texttt{R} 相同到 $[2,1]$；$[2,1]$
是刪除 \texttt{A}（$D[1,1] + 1 = 1$）；最後 \texttt{P} 相同到 $[0,0]$。 得到的正是上面的對齊。

\textbf{和 LCS 的關係。} 如果只允許插入和刪除，相同的字元配成一行，不同的字元只能一個刪、一
個插，成本 $2$，所以最好的對齊就是讓相同的行越多越好：
\[
\text{只用插入、刪除的距離} = m + n - 2 \cdot \text{LCS 長度} .
\]
\texttt{PARTNER} 和 \texttt{PRINTER} 的 LCS 長度是 $5$（例如 \texttt{PRTER}），只用插入、刪除
要 $7 + 7 - 10 = 4$ 步，允許替換時是 $3$ 步。

實際應用常把三種操作設成不同的成本（例如鍵盤上相鄰的字母替換得便宜一點），遞迴式只要把
三個 $1$ 換成對應的成本；再加上「相鄰兩個字元對調」也算一步的版本叫 Damerau 距離，要多看
$D[i-2, j-2]$。

\begin{example}[label=ex:dp-edit]{c}{(edit\_distance.c) 編輯距離與一種對齊}
#include <stdio.h>
#include <stdlib.h>
#include <string.h>

#define MAXLEN 1000

int min3(int a, int b, int c)
{
    int m = a < b ? a : b;
    return m < c ? m : c;
}

// 算出 x 改成 y 的編輯距離，並印出一種對齊（倒著存在 ax、ay、op）。
void edit(const char *x, const char *y)
{
    static char ax[2 * MAXLEN], ay[2 * MAXLEN], op[2 * MAXLEN];
    int m = strlen(x), n = strlen(y);
    int *d = malloc((size_t)(m + 1) * (n + 1) * sizeof *d);
    if (d == NULL) {
        return;
    }
#define D(i, j) d[(i) * (n + 1) + (j)]
    for (int i = 0; i <= m; i++) {
        D(i, 0) = i;                        // 刪掉 x 的前 i 個字
    }
    for (int j = 0; j <= n; j++) {
        D(0, j) = j;                        // 插入 y 的前 j 個字
    }
    for (int i = 1; i <= m; i++) {
        for (int j = 1; j <= n; j++) {
            int diff = (x[i - 1] != y[j - 1]);
            D(i, j) = min3(D(i - 1, j - 1) + diff, D(i - 1, j) + 1, D(i, j - 1) + 1);
        }
    }
    int i = m, j = n, k = 0;                // 依序試：左上、上（刪除）、左（插入）
    while (i > 0 || j > 0) {
        if (i > 0 && j > 0 && D(i, j) == D(i - 1, j - 1) + (x[i - 1] != y[j - 1])) {
            ax[k] = x[--i];
            ay[k] = y[--j];
            op[k] = (ax[k] == ay[k]) ? '=' : 'S';
        } else if (i > 0 && D(i, j) == D(i - 1, j) + 1) {
            ax[k] = x[--i];
            ay[k] = '-';
            op[k] = 'D';
        } else {
            ax[k] = '-';
            ay[k] = y[--j];
            op[k] = 'I';
        }
        k++;
    }
    printf("%s -> %s: distance %d\n", x, y, D(m, n));
#undef D
    free(d);
    const char *rows[3] = {ax, ay, op};
    for (int r = 0; r < 3; r++) {
        for (int t = k - 1; t >= 0; t--) {
            printf("%c%c", rows[r][t], t > 0 ? ' ' : '\n');
        }
    }
}

int main(void)
{
    static char x[MAXLEN + 1], y[MAXLEN + 1];
    while (scanf("%1000s %1000s", x, y) == 2) {
        edit(x, y);
    }
    return 0;
}
\end{example}
\iosample{PARTNER PRINTER}{PARTNER -> PRINTER: distance 3 \\ P A R T N - E R \\ P - R I N T E R \\ = D = S = I = =}

這組的最佳對齊只有一種：還原時每一步都只有一個方向符合。 有好幾種時，程式依照「左上、上、
左」的順序只印出其中一種。

\subsection{最長遞增子序列：\texorpdfstring{$O(n^2)$}{O(n²)}}
\label{sec:dp-lis}

給一個數列 $A[1..n]$，\textbf{最長遞增子序列}（longest increasing subsequence，LIS）是最長的
\textbf{嚴格遞增}子序列。 本節的例子是
\[
A = \langle 4, 9, 2, 7, 5, 3, 8, 6, 10, 1, 7 \rangle ,
\]
例如 $\langle 4, 7, 8, 10 \rangle$ 是一個長度 $4$ 的遞增子序列；這個數列沒有長度 $5$ 的。

只用前綴當子問題不夠：知道「$A[1..i]$ 的 LIS 長度」，還是不知道 $A[i+1]$ 接不接得上去，因為
接不接得上要看最後一個數多大。 所以子問題改成：令 $L[i]$ 為\textbf{以 $A[i]$ 結尾}的遞增子
序列的最大長度。 以 $A[i]$ 結尾的遞增子序列，去掉 $A[i]$ 後不是空的，就是以某個 $A[j]$ 結尾
（$j < i$、$A[j] < A[i]$）的遞增子序列，而且那一段要盡量長：
\[
L[i] = 1 + \max\bigl\{\, L[j] : 1 \le j < i,\ A[j] < A[i] \,\bigr\},
\]
大括號裡沒有任何 $j$ 時，最大值當作 $0$。 答案是 $\max_i L[i]$，不是 $L[n]$。

\begin{pseudo}[label=alg:dp-lis-quad]{LIS，$O(n^2)$}
\Function{LIS-Quadratic}{$A, n$}
    \For{$i \gets 1$ \To $n$}
        \State $L[i] \gets 1$
        \State $\mathit{prev}[i] \gets \NIL$
        \For{$j \gets 1$ \To $i - 1$}
            \If{$A[j] < A[i]$ \AND $L[j] + 1 > L[i]$}
                \State $L[i] \gets L[j] + 1$
                \State $\mathit{prev}[i] \gets j$
            \EndIf
        \EndFor
    \EndFor
    \State \Return $\max_{1 \le i \le n} L[i]$
\EndFunction
\end{pseudo}

$\mathit{prev}[i]$ 記下 $A[i]$ 前面接的是誰，平手時保留最小的 $j$（第 6 行用 $>$）。 兩層迴圈，
時間 $\Theta(n^2)$。

\begin{center}
\begin{tabular}{l*{11}{c}}
\toprule
$i$ & 1 & 2 & 3 & 4 & 5 & 6 & 7 & 8 & 9 & 10 & 11 \\
\midrule
$A[i]$ & 4 & 9 & 2 & 7 & 5 & 3 & 8 & 6 & 10 & 1 & 7 \\
$L[i]$ & 1 & 2 & 1 & 2 & 2 & 2 & 3 & 3 & 4 & 1 & 4 \\
$\mathit{prev}[i]$ & -- & 1 & -- & 1 & 1 & 3 & 4 & 5 & 7 & -- & 8 \\
\bottomrule
\end{tabular}
\end{center}

例如 $L[7]$：$A[7] = 8$，前面比 $8$ 小的有 $A[1] = 4$、$A[3] = 2$、$A[4] = 7$、$A[5] = 5$、
$A[6] = 3$，它們的 $L$ 是 $1, 1, 2, 2, 2$，最大的 $2$ 最先出現在 $j = 4$，所以 $L[7] = 3$、
$\mathit{prev}[7] = 4$。 最大的 $L$ 是 $L[9] = 4$（$L[11]$ 也是 $4$，取先出現的），沿著
$\mathit{prev}$ 走回去：$9 \to 7 \to 4 \to 1$，LIS 是 $\langle 4, 7, 8, 10 \rangle$。

另一個看法：LIS 就是 $A$ 和「把 $A$ 排序、去掉重複」之後的序列的 LCS。 這也是 $O(n^2)$。

\subsection{最長遞增子序列：\texorpdfstring{$O(n \lg n)$}{O(n lg n)}}
\label{sec:dp-lis-fast}

$n$ 到 $10^6$ 時 $n^2$ 太慢。 換一個子問題：由左往右掃，對每個長度 $k$，記住
\[
\mathit{tail}[k] = \text{目前看過的數裡，所有長度 $k$ 的遞增子序列中，結尾最小的那一個的位置。}
\]
同樣長度時，結尾越小越好，因為之後越容易接上新的數，所以只記結尾最小的那一個就夠了。

\begin{whybox}[{tail 的值一定是嚴格遞增的}]
設目前有長度 $k$ 和長度 $k+1$ 的遞增子序列。 取一個長度 $k + 1$、結尾最小的，把它的最後一個
數拿掉，得到一個長度 $k$ 的遞增子序列，它的結尾比原本的結尾小。 所以長度 $k$ 的最小結尾，嚴
格小於長度 $k + 1$ 的最小結尾：
\[
A[\mathit{tail}[1]] < A[\mathit{tail}[2]] < \cdots < A[\mathit{tail}[\mathit{len}]] ,
\]
$\mathit{len}$ 是目前最長的長度。 這個陣列是排好序的，所以可以二分搜尋。

處理新的數 $x = A[i]$ 時，找最小的 $k$ 使 $A[\mathit{tail}[k]] \ge x$：
\begin{itemize}
\item $x$ 可以接在長度 $k - 1$、結尾 $A[\mathit{tail}[k-1]] < x$ 的那一條後面，得到長度 $k$、
結尾 $x$ 的遞增子序列，而且 $x \le A[\mathit{tail}[k]]$，所以 $x$ 成為新的長度 $k$ 最小結尾。
\item 長度小於 $k$ 的最小結尾都比 $x$ 小，$x$ 換不掉它們。 長度 $k$ 以上的結尾都 $\ge x$，$x$
接不上去，所以也生不出更長、以 $x$ 結尾的遞增子序列。 其他長度的最小結尾都不變。
\item 都比 $x$ 小時 $k = \mathit{len} + 1$，最長的長度加一。
\end{itemize}
每一步之後 $\mathit{tail}$ 的定義仍然成立，所以掃完時 $\mathit{len}$ 就是 LIS 的長度。
\end{whybox}

「最小的 $k$ 使 $A[\mathit{tail}[k]] \ge x$」正是第~\ref{sec:dc-bsearch}~節 \proc{Lower-Bound}
做的事，每次 $O(\lg n)$，總時間 $O(n \lg n)$。 為了還原，另外記 $\mathit{prev}[i]$：$A[i]$ 放進
第 $k$ 格時，它前面接的是當時的 $\mathit{tail}[k - 1]$。

\begin{pseudo}[label=alg:dp-lis-fast]{LIS，$O(n \lg n)$}
\Function{LIS-Fast}{$A, n$}
    \State $\mathit{len} \gets 0$
    \For{$i \gets 1$ \To $n$}
        \State 二分搜尋最小的 $k$（$1 \le k \le \mathit{len}$）使 $A[\mathit{tail}[k]] \ge A[i]$；沒有就令 $k \gets \mathit{len} + 1$
        \If{$k > 1$}
            \State $\mathit{prev}[i] \gets \mathit{tail}[k - 1]$
        \Else
            \State $\mathit{prev}[i] \gets \NIL$
        \EndIf
        \State $\mathit{tail}[k] \gets i$
        \If{$k = \mathit{len} + 1$}
            \State $\mathit{len} \gets \mathit{len} + 1$
        \EndIf
    \EndFor
    \State \Return $\mathit{len}$
\EndFunction
\end{pseudo}

同樣的數列跑一次。 表中「tail 的值」一欄寫的是 $A[\mathit{tail}[1]], A[\mathit{tail}[2]],
\ldots$，方便看；程式裡存的是位置。

\begin{center}
\begin{tabular}{cclclc}
\toprule
$i$ & $A[i]$ & $k$ & 動作 & tail 的值 & $\mathit{prev}[i]$ \\
\midrule
1 & 4 & 1 & 接在最後 & 4 & -- \\
2 & 9 & 2 & 接在最後 & 4, 9 & 1 \\
3 & 2 & 1 & 取代第 1 格 & 2, 9 & -- \\
4 & 7 & 2 & 取代第 2 格 & 2, 7 & 3 \\
5 & 5 & 2 & 取代第 2 格 & 2, 5 & 3 \\
6 & 3 & 2 & 取代第 2 格 & 2, 3 & 3 \\
7 & 8 & 3 & 接在最後 & 2, 3, 8 & 6 \\
8 & 6 & 3 & 取代第 3 格 & 2, 3, 6 & 6 \\
9 & 10 & 4 & 接在最後 & 2, 3, 6, 10 & 8 \\
10 & 1 & 1 & 取代第 1 格 & 1, 3, 6, 10 & -- \\
11 & 7 & 4 & 取代第 4 格 & 1, 3, 6, 7 & 8 \\
\bottomrule
\end{tabular}
\end{center}

最後 $\mathit{len} = 4$，$\mathit{tail}[4] = 11$。 沿 $\mathit{prev}$ 走回去是位置
$11, 8, 6, 3$，倒過來讀，LIS 是 $A[3]$、$A[6]$、$A[8]$、$A[11]$，也就是 $\langle 2, 3, 6, 7 \rangle$。
它和 $O(n^2)$ 版本找到的 $\langle 4, 7, 8, 10 \rangle$ 不同，但一樣長；遇到好幾個 LIS 時，兩種
方法選的不一樣。

要注意兩件事。 第一，最後的 tail 值 $1, 3, 6, 7$ \textbf{不是}一個子序列：$1$ 在原數列裡排在
$3$、$6$ 後面。 tail 只保證長度對，要還原必須用 $\mathit{prev}$。 第二，相等的數怎麼處理：嚴格遞增時相等的數不能接在一起，所以找的是第一個「$\ge$」的格子，
遇到和某個 tail 值相等的數時，是\textbf{取代}那一格，不會接到它後面。 如果題目要的是
\textbf{非嚴格}遞增（允許相等），要改找第一個「$>$」的格子。

\subsection{C 程式}

\texttt{tail} 存的是位置，比較時用 \texttt{a[tail[mid]]}。 二分搜尋的寫法和範例~\ref{ex:dc-bsearch}
的 \texttt{lower\_bound} 一樣，範圍是左閉右開的 $[0, \mathit{len})$。

\begin{example}[label=ex:dp-lis]{c}{(lis.c) $O(n \lg n)$ 的最長遞增子序列}
#include <stdio.h>
#include <stdlib.h>

#define MAXN 1000000

// 最長嚴格遞增子序列，O(n lg n)。out 放一個 LIS，回傳長度。
// tail[k]：長度 k+1 的遞增子序列裡結尾最小的那一個，結尾的位置。
// pred[i]：以 a[i] 結尾的那一條，前一個數的位置（-1 表示沒有）。
int lis(const int *a, int n, int *out)
{
    int *tail = malloc(n * sizeof *tail);
    int *pred = malloc(n * sizeof *pred);
    if (tail == NULL || pred == NULL) {
        free(tail);
        free(pred);
        return -1;
    }
    int len = 0;
    for (int i = 0; i < n; i++) {
        int lo = 0, hi = len;               // 找第一個 a[tail[k]] >= a[i] 的 k
        while (lo < hi) {
            int mid = lo + (hi - lo) / 2;
            if (a[tail[mid]] < a[i]) {
                lo = mid + 1;
            } else {
                hi = mid;
            }
        }
        pred[i] = (lo > 0) ? tail[lo - 1] : -1;
        tail[lo] = i;                       // 取代，或接在最後面
        if (lo == len) {
            len++;
        }
    }
    for (int k = len - 1, i = tail[len - 1]; k >= 0; k--, i = pred[i]) {
        out[k] = a[i];                      // 沿著 pred 倒著走回去
    }
    free(tail);
    free(pred);
    return len;
}

int main(void)
{
    static int a[MAXN], out[MAXN];
    int n;
    while (scanf("%d", &n) == 1 && n > 0 && n <= MAXN) {
        for (int i = 0; i < n; i++) {
            if (scanf("%d", &a[i]) != 1) {
                return 1;
            }
        }
        int len = lis(a, n, out);
        printf("length %d:", len);
        for (int k = 0; k < len; k++) {
            printf(" %d", out[k]);
        }
        printf("\n");
    }
    return 0;
}
\end{example}
\iosample{11 \\ 4 9 2 7 5 3 8 6 10 1 7 \\ 6 \\ 5 5 5 5 5 5 \\ 7 \\ 9 8 7 6 5 4 3}{length 4: 2 3 6 7 \\ length 1: 5 \\ length 1: 3}

第二組六個 $5$：嚴格遞增時相等的數不能接，LIS 長度是 $1$。 第三組完全遞減，每個數都取代第
一格，長度也是 $1$。

\section{所有點對最短路徑的 DP 觀點}
\label{sec:dp-apsp}

最後用這一章的眼光看一個圖論問題。 有向圖的點編號 $1$ 到 $n$，邊 $(i, j)$ 有權重 $w_{ij}$（可以
是負的，但圖裡沒有總權重為負的環）。 \textbf{所有點對最短路徑}問題要算出每一對 $(i, j)$ 的最
短路徑長度。 第~25~章會完整講這個問題，包括怎麼還原路徑、怎麼偵測負環，以及其他做法；這
一節只看它的遞迴式從哪裡來。

\subsection{子問題：限制中間點}

一條路徑除了起點和終點以外經過的點，叫它的\textbf{中間點}。 定義
\[
d_{ij}^{(k)} = \text{從 } i \text{ 到 } j \text{、中間點都在 } \{1, 2, \ldots, k\} \text{ 裡的最短路徑長度。}
\]
$k = 0$ 時不准有中間點，只能直接走一條邊：$d_{ij}^{(0)} = w_{ij}$，$i = j$ 時是 $0$，沒有邊時
是 $\infty$。 $k = n$ 時沒有限制，$d_{ij}^{(n)}$ 就是答案。

從 $k - 1$ 到 $k$：看中間點都在 $\{1, \ldots, k\}$ 的最短路徑 $p$。 沒有負環時可以假設 $p$
是簡單路徑，所以點 $k$ 在 $p$ 上最多出現一次。
\begin{itemize}
\item $p$ 不經過 $k$：它的中間點都在 $\{1, \ldots, k-1\}$，長度是 $d_{ij}^{(k-1)}$。
\item $p$ 經過 $k$：在 $k$ 切成 $i \leadsto k$ 和 $k \leadsto j$ 兩段，兩段的中間點都在
$\{1, \ldots, k - 1\}$（$k$ 只出現一次，現在是端點），而且各自都要是最短的（剪下貼上），長度是
$d_{ik}^{(k-1)} + d_{kj}^{(k-1)}$。
\end{itemize}
所以
\[
d_{ij}^{(k)} = \min\bigl( d_{ij}^{(k-1)},\ d_{ik}^{(k-1)} + d_{kj}^{(k-1)} \bigr) .
\]
這和 0/1 背包的「第 $k$ 件拿或不拿」是同一種想法：第 $k$ 個點「用或不用」。 和第~\ref{sec:dp-substructure}~節最
長簡單路徑的反例對照：最短路徑的兩段即使經過同一個點，也能把繞的那一圈拿掉而不變長，所以
子問題是獨立的；這正是「沒有負環」這個條件的作用。

子問題有 $n^3$ 個，每個 $O(1)$，時間 $\Theta(n^3)$。 第 $k$ 層只用到第 $k - 1$ 層，而且第 $k$
層的第 $k$ 列、第 $k$ 欄和第 $k - 1$ 層相同：以第 $k$ 欄為例，$d_{kk}^{(k-1)} = 0$，所以
$d_{ik}^{(k)} = \min\bigl(d_{ik}^{(k-1)},\ d_{ik}^{(k-1)} + 0\bigr) = d_{ik}^{(k-1)}$，第 $k$ 列同理。
第 $k$ 層要讀的只有第 $k$ 列、第 $k$ 欄，它們在這一層不會變，所以可以只用一個矩陣就地更新，
空間 $\Theta(n^2)$。 這就是 \textbf{Floyd--Warshall 演算法}：

\begin{pseudo}[label=alg:dp-fw]{Floyd--Warshall（就地更新）}
\Function{Floyd-Warshall}{$W, n$}
    \State $D \gets W$ \Comment{複製一份 $n \times n$ 矩陣}
    \For{$k \gets 1$ \To $n$}
        \For{$i \gets 1$ \To $n$}
            \For{$j \gets 1$ \To $n$}
                \If{$D[i, k] + D[k, j] < D[i, j]$}
                    \State $D[i, j] \gets D[i, k] + D[k, j]$
                \EndIf
            \EndFor
        \EndFor
    \EndFor
    \State \Return $D$
\EndFunction
\end{pseudo}

最外層一定要是 $k$。 把 $k$ 放到最內層，算 $D[i, j]$ 時用到的 $D[i, k]$、$D[k, j]$ 還沒有考慮
過足夠多的中間點，答案會錯，這是寫這個演算法最常見的錯誤。

\subsection{一個例子}

四個點，邊是 $1 \to 2$（$4$）、$1 \to 3$（$9$）、$2 \to 3$（$3$）、$2 \to 4$（$8$）、
$3 \to 1$（$6$）、$3 \to 4$（$-2$）、$4 \to 1$（$5$）。 依序算出 $D^{(0)}$ 到 $D^{(4)}$，粗體是
這一層變小的格子：
\begin{gather*}
D^{(0)} = \begin{pmatrix} 0 & 4 & 9 & \infty \\ \infty & 0 & 3 & 8 \\ 6 & \infty & 0 & -2 \\ 5 & \infty & \infty & 0 \end{pmatrix},\quad D^{(1)} = \begin{pmatrix} 0 & 4 & 9 & \infty \\ \infty & 0 & 3 & 8 \\ 6 & \mathbf{10} & 0 & -2 \\ 5 & \mathbf{9} & \mathbf{14} & 0 \end{pmatrix},\quad D^{(2)} = \begin{pmatrix} 0 & 4 & \mathbf{7} & \mathbf{12} \\ \infty & 0 & 3 & 8 \\ 6 & 10 & 0 & -2 \\ 5 & 9 & \mathbf{12} & 0 \end{pmatrix}, \\
D^{(3)} = \begin{pmatrix} 0 & 4 & 7 & \mathbf{5} \\ \mathbf{9} & 0 & 3 & \mathbf{1} \\ 6 & 10 & 0 & -2 \\ 5 & 9 & 12 & 0 \end{pmatrix},\quad D^{(4)} = \begin{pmatrix} 0 & 4 & 7 & 5 \\ \mathbf{6} & 0 & 3 & 1 \\ \mathbf{3} & \mathbf{7} & 0 & -2 \\ 5 & 9 & 12 & 0 \end{pmatrix}
\end{gather*}

例如 $D^{(3)}$ 的 $d_{14}$ 從 $12$（$1 \to 2 \to 4$）變成 $5$：$d_{13}^{(2)} + d_{34}^{(2)} =
7 + (-2)$，也就是 $1 \to 2 \to 3 \to 4$。 $D^{(4)}$ 的 $d_{31} = 3$ 是 $3 \to 4 \to 1$，比直接
的邊 $6$ 短。 每一層都有格子變小，最後一層就是答案。

子問題也可以改成「最多用 $k$ 條邊」（這裡的 $k$ 是邊數，不是中間點的編號）：
$\ell_{ij}^{(k)} = \min_{t} \bigl(\ell_{it}^{(k-1)} + w_{tj}\bigr)$，看最後一條邊從哪個點 $t$
走進 $j$。 要算到 $k = n - 1$，每層 $\Theta(n^3)$，共 $\Theta(n^4)$；它的形狀和矩陣乘法一樣，
用第~\ref{sec:dc-power}~節快速冪的反覆平方可以降到 $\Theta(n^3 \lg n)$。 同一個問題，
子問題選得不同，時間可以差很多，第~25~章會比較這幾種做法。

\subsection{本章的 DP 整理}

\begin{center}
\small
\begin{tabular}{>{\raggedright\arraybackslash}p{0.2\linewidth}>{\raggedright\arraybackslash}p{0.22\linewidth}cc>{\raggedright\arraybackslash}p{0.22\linewidth}}
\toprule
問題 & 子問題 & 子問題數 & 每個的選擇 & 時間 \\
\midrule
二項式係數 & $\binom{i}{j}$ & $\Theta(nk)$ & $1$ & $\Theta(nk)$ \\
矩陣鏈乘 & 區間 $A_{i..j}$ & $\Theta(n^2)$ & $O(n)$ & $\Theta(n^3)$ \\
LCS & 兩個前綴 $X_i, Y_j$ & $\Theta(mn)$ & $O(1)$ & $\Theta(mn)$ \\
最佳二元搜尋樹 & 區間 $k_i..k_j$ & $\Theta(n^2)$ & $O(n)$ & $\Theta(n^3)$，Knuth 改進後 $\Theta(n^2)$ \\
0/1 背包 & 前 $i$ 件、容量 $c$ & $\Theta(nW)$ & $2$ & $\Theta(nW)$，擬多項式 \\
編輯距離 & 兩個前綴 $X_i, Y_j$ & $\Theta(mn)$ & $3$ & $\Theta(mn)$ \\
LIS & 以 $A[i]$ 結尾 & $n$ & $O(n)$ & $\Theta(n^2)$；改寫後 $O(n \lg n)$ \\
Floyd--Warshall & 中間點限 $\{1..k\}$ & $\Theta(n^3)$ & $2$ & $\Theta(n^3)$ \\
\bottomrule
\end{tabular}
\end{center}

常見的子問題只有幾種形狀：\textbf{前綴}（LCS、編輯距離、背包的前 $i$ 件）、\textbf{區間}（矩陣
鏈乘、最佳二元搜尋樹）、\textbf{以某個元素結尾}（LIS），以及多加一個\textbf{數值參數}（背包的容
量、Floyd--Warshall 的 $k$）。 遇到新問題時，先試這幾種，再看最佳解的「最後一步」有幾種選
擇。 只要值時，前綴型和背包型的表通常都可以只留一兩列；區間型的表依賴的格子散在各列，不容
易這樣省。

\section{本章常見錯誤}
\label{sec:dp-pitfalls}

\begin{longtable}{p{0.25\linewidth}p{0.48\linewidth}>{\raggedright\arraybackslash}p{0.17\linewidth}}
\toprule
錯誤 & 例子 & 見 \\
\midrule
\endhead
沒有定義表格的一格 & 只寫遞迴式，沒寫 $c[i, j]$ 代表什麼，或邊界漏寫 &
第~\ref{sec:dp-steps}~節 \\
記憶化的記號和答案撞在一起 & 用 $0$ 表示「還沒算過」，但子問題的答案本來就可能是 $0$，於
是一直重算 & 第~\ref{sec:dp-memo}~節 \\
填表順序錯 & 矩陣鏈乘一列一列由上往下填，算 $m[1, 3]$ 時 $m[2, 3]$ 還沒算 &
第~\ref{sec:dp-matrix}~節 \\
把一個候選當成最佳 & 矩陣鏈乘一律在最小的維度切開；最佳二元搜尋樹一律讓最常查的鍵當根 &
第~\ref{sec:dp-matrix}、\ref{sec:dp-obst}~節 \\
還原 LCS 的規則不對 & 沒先檢查 $x_i = y_j$ 就比較上、左兩格，漏掉該選的字元；或沒照題目規定的平手方向，還原出另一個 LCS &
第~\ref{sec:dp-lcs}~節 \\
最佳二元搜尋樹的成本模型混用 & 外部節點用深度加一算成本，空樹卻寫 $e[i, i-1] = 0$；
$w(i, j)$ 的 $q$ 從 $q_i$ 開始加 & 第~\ref{sec:dp-obst}~節 \\
一維背包的方向 & 0/1 背包的容量由小到大跑，同一件物品被拿很多次 &
第~\ref{sec:dp-knapsack-1d}~節 \\
說背包 DP 是多項式時間 & $\Theta(nW)$ 是 $W$ 的數值的多項式，不是位元數的 &
第~\ref{sec:dp-pseudo-poly}~節 \\
編輯距離的邊界 & 把 $D[i, 0]$、$D[0, j]$ 都設成 $0$ & 第~\ref{sec:dp-edit}~節 \\
LIS 的答案 & 答案寫成 $L[n]$，不是 $\max_i L[i]$；把最後的 tail 陣列當成一個 LIS &
第~\ref{sec:dp-lis}、\ref{sec:dp-lis-fast}~節 \\
嚴格和非嚴格遞增 & 嚴格遞增卻找第一個 $>$ 的格子，相等的數被接在一起 &
第~\ref{sec:dp-lis-fast}~節 \\
Floyd--Warshall 的迴圈順序 & 把 $k$ 放在最內層 & 第~\ref{sec:dp-apsp}~節 \\
\bottomrule
\end{longtable}

\section{觀念題}
\label{sec:dp-exercises}

題型和研究所考試、期中期末考的筆試相近，有選擇、計算、追蹤、設計和證明。 平手時一律用本章
的規則。 解答在本章最後。

\begin{enumerate}
\item \label{dp-q-tf}（選擇，可能不只一個答案）下列哪些敘述正確？
\begin{enumerate}
\item 同一個遞迴式，記憶化和由下而上填表的漸近時間一定相同。
\item 0/1 背包的 $\Theta(nW)$ 演算法是多項式時間演算法。
\item 矩陣鏈乘只要求 $m[1, n]$ 時，可以像 LCS 一樣只留兩列。
\item Floyd--Warshall 的三層迴圈裡，$k$ 必須在最外層。
\item 一個最佳化問題只要有最佳子結構，就能用動態規劃在多項式時間內解出來。
\item 把 $X$、$Y$ 兩個字串對調，LCS 的長度不變，還原出來的 LCS 也不變。
\end{enumerate}

\item \label{dp-q-mc}（計算）四個矩陣的大小是 $\langle 5, 8, 2, 10, 4 \rangle$。 依鏈長填出
$m$ 表和 $s$ 表（寫出每一格每個 $k$ 的成本），寫出最佳的括號方式，並和由左往右乘
$((A_1 A_2) A_3) A_4$ 的次數比較。

\item \label{dp-q-mc-rec}（證明）不用記憶化，直接照 $m[i, j]$ 的遞迴式寫遞迴函式算
$m[1, n]$。 令 $R(n)$ 為鏈長 $n$ 時函式被呼叫的總次數（$R(1) = 1$）。 寫出 $R(n)$ 的遞迴
式，證明 $R(n) = 3^{n-1}$，並和記憶化時不同子問題的個數比較。

\item \label{dp-q-lcs}（追蹤）$X = \texttt{TEACHER}$（列）、$Y = \texttt{CHEATER}$（欄）。
\begin{enumerate}
\item 用 \proc{LCS-Length} 填出 $c$ 表和箭頭，用 \proc{Print-LCS} 還原一個 LCS。
\item 列出所有的 LCS。
\item 平手時改成指向左，還原出來的是哪一個？
\end{enumerate}

\item \label{dp-q-obst}（計算）三個鍵 $k_1 < k_2 < k_3$，查詢次數
$p = \langle 4, 2, 6 \rangle$，查不到的次數 $q = \langle 1, 3, 2, 2 \rangle$（$q_0$ 到 $q_3$）。
用本章的成本模型填出 $w$、$e$、$\mathit{root}$ 表，畫出最佳二元搜尋樹（含外部節點），算出
平均每次查詢的比較次數。 如果改用「走到外部節點也算一次」的成本模型，$e[1, 3]$ 是多少？

\item \label{dp-q-ks}（追蹤）背包容量 $W = 8$，四件物品的 $(w_i, v_i)$ 是 $(2, 5)$、
$(5, 9)$、$(4, 10)$、$(3, 10)$。
\begin{enumerate}
\item 填出 $K$ 表，用 \proc{Knapsack-Items} 找出要拿的物品。
\item 用一維陣列但容量由小到大跑，$d[0..8]$ 最後是什麼？ 這是哪個問題的答案？
\end{enumerate}

\item \label{dp-q-sparse}（計算）$n = 20$ 件物品，容量 $W = 10^9$，每件的重量都在
$10^7$ 到 $10^8$ 之間。
\begin{enumerate}
\item \proc{Knapsack} 的表有幾格？ 開得出來嗎？
\item 改用記憶化：$\proc{Best}(i, c)$ 回傳「前 $i$ 件、容量 $c$」的最大價值，從
$\proc{Best}(20, W)$ 開始遞迴。 證明會被呼叫到的不同 $(i, c)$ 最多 $2^{21} - 1$ 個。 這時表
要用什麼資料結構存？
\end{enumerate}

\item \label{dp-q-edit}（追蹤）填出 \texttt{SPORT}（列）改成 \texttt{STRONG}（欄）的 $D$ 表，
寫出編輯距離和一種最少步數的改法（用上下兩列的對齊表示）。

\item \label{dp-q-edit-match}（證明）證明 $x_i = y_j$ 時 $D[i, j] = D[i-1, j-1]$。 提示：先
證明 $D[i-1, j-1] \le D[i-1, j] + 1$ 和 $D[i-1, j-1] \le D[i, j-1] + 1$。

\item \label{dp-q-lis}（追蹤）$A = \langle 3, 8, 4, 12, 6, 10, 7, 11, 5 \rangle$。
\begin{enumerate}
\item 用 \proc{LIS-Quadratic} 算出 $L$ 和 $\mathit{prev}$，還原一個 LIS。
\item 用 \proc{LIS-Fast} 寫出每一步之後的 tail 值和 $\mathit{prev}$，還原一個 LIS。
\item 最後的 tail 值是不是 $A$ 的子序列？
\end{enumerate}

\item \label{dp-q-alloc}（設計）期末考前還有 $T = 4$ 小時，要分給三科，每科只能分整數小
時。 第 $i$ 科讀 $h$ 小時預計拿到的分數 $f_i(h)$ 如下表。 定義子問題、寫出遞迴式，填表求出
最高總分和每科分到的時數，並寫出一般 $n$ 科、$T$ 小時的時間複雜度。
\begin{center}
\begin{tabular}{c*{5}{c}}
\toprule
$h$ & 0 & 1 & 2 & 3 & 4 \\
\midrule
$f_1(h)$ & 0 & 5 & 8 & 10 & 11 \\
$f_2(h)$ & 0 & 3 & 7 & 10 & 12 \\
$f_3(h)$ & 0 & 4 & 6 & 7 & 8 \\
\bottomrule
\end{tabular}
\end{center}

\item \label{dp-q-fw}（追蹤）三個點，邊是 $1 \to 2$（$5$）、$1 \to 3$（$4$）、$2 \to 1$
（$2$）、$2 \to 3$（$-3$）、$3 \to 1$（$4$）。 寫出 $D^{(0)}$、$D^{(1)}$、$D^{(2)}$、$D^{(3)}$。

\item \label{dp-q-sub}（選擇並說明）下列哪些問題有最佳子結構？
\begin{enumerate}
\item 有向無環圖上，從 $s$ 到 $t$ 邊數最多的路徑。
\item 一般有向圖上，從 $s$ 到 $t$ 邊數最多的簡單路徑。
\item 帶權有向圖上，從 $s$ 到 $t$、恰好走 $k$ 條邊（點可以重複）的最小權重走法。
\end{enumerate}
\end{enumerate}

\section{實作練習}
\label{sec:dp-practice}

下面是 Aki-Judge 上和本章相關的題目。 矩陣鏈乘和最佳二元搜尋樹目前沒有對應的練習題，請用觀
念題第~\ref{dp-q-mc}、\ref{dp-q-obst}~題練習。

\begin{center}
\begin{tabular}{llc}
\toprule
小節 & 題目 & 難度 \\
\midrule
DP 的要素 & \akiref{cheapest-route-across-a-grid}{網格上最便宜的路線} & 暖身 \\
DP 的要素 & \akiref{cheapest-grid-route-with-blocked-cells}{含封鎖格的最便宜路線} & 基本 \\
DP 的要素 & \akiref{mandatory-checkpoints-on-a-down-right-route}{送貨機器人的必經格} & 進階 \\
DP 的要素 & \akiref{heaviest-set-of-non-adjacent-tree-nodes}{樹上互不相鄰節點的最大權重和} & 進階 \\
LCS & \akiref{length-of-the-longest-common-subsequence}{最長共同子序列長度} & 基本 \\
LCS & \akiref{common-subsequence-without-repeated-neighbors}{相鄰不重複的共同子序列} & 進階 \\
0/1 背包 & \akiref{best-value-within-a-weight-limit}{重量上限內的最大價值} & 基本 \\
0/1 背包 & \akiref{best-value-with-unlimited-copies}{可重複取用的最大價值} & 基本 \\
編輯距離 & \akiref{minimum-edits-between-two-strings}{兩字串之間的最少編輯次數} & 基本 \\
LIS & \akiref{longest-strictly-increasing-subsequence-small-n}{最長嚴格遞增子序列（小範圍）} & 基本 \\
LIS & \akiref{longest-strictly-increasing-subsequence-large-n}{最長嚴格遞增子序列（大範圍）} & 進階 \\
\bottomrule
\end{tabular}
\end{center}

幾點提示：
\begin{itemize}
\item 兩題網格路線：每一步只能往右或往下，所以一格只會從上面或左邊走進來，照列、照欄的順
序填就好。 封鎖格當成成本 $\infty$，走不到終點時輸出題目指定的值。 「必經格」那一題要兩張
表：從起點走不走得到這一格、從這一格走不走得到終點，兩者都成立的格子才在某條路線上。
\item 樹上的題目，子問題是「以 $v$ 為根的子樹，$v$ 選或不選」兩種狀態。 $n$ 到 $2 \times 10^5$，
樹可能是一條長鏈，遞迴太深，請先用 BFS 排出順序，再倒著由葉往根算。
\item 「相鄰不重複的共同子序列」是 LCS 的變形：除了兩個位置，子問題還要記住上一個選進來的
字母（$26$ 種加上「還沒選」），用記憶化最順手。
\item 0/1 背包那題 $n \le 1000$、$W \le 10^5$，二維表太大，要用第~\ref{sec:dp-knapsack-1d}~節的一維
陣列、容量由大到小；可重複取用那題正好相反。
\item 兩題 LCS、編輯距離的字串長度到 $2000$，整張表約四百萬格，可以開；只留兩列更省。
\item LIS 大範圍那題 $n$ 到 $10^6$，要用 $O(n \lg n)$ 的做法，注意是嚴格遞增。
\end{itemize}

\section{讀完你應該會}
\label{sec:dp-checklist}

\begin{itemize}
\item 判斷一個問題有沒有重疊的子問題和最佳子結構，用剪下貼上證明最佳子結構，並舉出最長簡
單路徑這種沒有最佳子結構的例子。
\item 寫動態規劃的答案時，用一句話定義表格的一格，寫出遞迴式、邊界、填表順序、答案位置和
複雜度。
\item 對矩陣鏈乘和最佳二元搜尋樹依區間長度填表，由 $s$ 或 $\mathit{root}$ 表還原括號或樹，
並說清楚最佳二元搜尋樹用的成本模型。
\item 填 LCS 和編輯距離的表，照指定的平手規則還原一個解；只要值時改成只留兩列。
\item 寫出 0/1 背包的二維表和一維版本，說明一維時容量為什麼由大到小，以及 $\Theta(nW)$ 為
什麼是擬多項式時間。
\item 寫出 LIS 的 $O(n^2)$ 和 $O(n \lg n)$ 做法並還原一個 LIS；寫出 Floyd--Warshall 的遞迴式，
說明 $k$ 為什麼在最外層。
\end{itemize}

\textbf{延伸閱讀}（書末會整理完整的參考資料）

\begin{itemize}
\item Thomas H. Cormen 等，\textit{Introduction to Algorithms}，第 4 版，第 14 章：動態規劃的
四個經典例子和「要素」一節的完整討論；所有點對最短路徑在第 23 章。 它的最佳二元搜尋樹把走
到外部節點也算一次比較，和本章差一個常數，對照時要注意。
\item Steven S. Skiena，\textit{The Algorithm Design Manual}，第 3 版，第 10 章：用很多實際例
子講怎麼「想出」子問題，包括編輯距離的各種變形。
\item Antti Laaksonen，\textit{Competitive Programmer's Handbook}，第 7 章：篇幅短，有 LIS、
網格路徑、背包和編輯距離的程式寫法。
\item D. E. Knuth，``Optimum Binary Search Trees''，\textit{Acta Informatica} 1(1)，1971：最佳
二元搜尋樹的 $\Theta(n^2)$ 方法，也就是本章提到的 $\mathit{root}$ 單調性。
\item D. S. Hirschberg，``A Linear Space Algorithm for Computing Maximal Common
Subsequences''，\textit{Communications of the ACM} 18(6)，1975：只用線性空間就能還原 LCS 的
分治方法。
\end{itemize}

\addcontentsline{toc}{section}{觀念題解答}
\section*{觀念題解答}

\begin{enumerate}
\item (d) 正確。
(a) 錯：記憶化只算真的用得到的子問題，可能比整張表少很多（第~\ref{dp-q-sparse}~題）；反過
來說，兩者的上界相同。 (b) 錯：是擬多項式。 (c) 錯：$m[i, j]$ 要用到第 $i$ 列左邊的所有
格子和第 $j$ 欄下面的所有格子，散在很多列，只留兩列不夠。 (d) 對：第 $k$ 輪要用到的
$D[i, k]$、$D[k, j]$ 必須已經是「中間點在 $\{1..k-1\}$」的最短值。 (e) 錯：最佳子結構只保證
遞迴式寫得出來，子問題可能有指數多個，例如旅行推銷員問題（第~20~章：每個城市走一次再回到起點的最短路
線）用動態規劃時，子問題要記住「走過的點的集合」，有 $2^n$ 種。
(f) 錯：長度不變，但平手規則的方向跟著對調，還原出來的可能不同，範例~\ref{ex:dp-lcs} 的前兩
組就是例子。

\item 鏈長 $2$：$m[1,2] = 5 \cdot 8 \cdot 2 = 80$、$m[2,3] = 8 \cdot 2 \cdot 10 = 160$、
$m[3,4] = 2 \cdot 10 \cdot 4 = 80$。
鏈長 $3$：$m[1,3]$：$k = 1$ 得 $0 + 160 + 400 = 560$，$k = 2$ 得 $80 + 0 + 100 = 180$，所以
$180$、$s = 2$。 $m[2,4]$：$k = 2$ 得 $0 + 80 + 64 = 144$，$k = 3$ 得 $160 + 0 + 320 = 480$，所
以 $144$、$s = 2$。
鏈長 $4$：$m[1,4]$：$k = 1$ 得 $0 + 144 + 160 = 304$，$k = 2$ 得 $80 + 80 + 40 = 200$，
$k = 3$ 得 $180 + 0 + 200 = 380$，所以 $200$、$s = 2$。
最佳是 $((A_1 A_2)(A_3 A_4))$，$200$ 次。 由左往右乘正是 $k = 3$ 那個候選，$380$ 次。

\item 每次呼叫自己算一次，再對每個 $k$ 呼叫左右兩邊：
\[
R(n) = 1 + \sum_{k=1}^{n-1} \bigl(R(k) + R(n - k)\bigr) = 1 + 2\sum_{k=1}^{n-1} R(k) .
\]
$n \ge 2$ 時把 $R(n)$ 和 $R(n-1)$ 的式子相減（$R(1) = 1$ 也符合 $R(1) = 1 + 2 \cdot 0$），得
$R(n) - R(n-1) = 2R(n-1)$，即 $R(n) = 3R(n-1)$，所以 $R(n) = 3^{n-1}$。 驗算：
$R(2) = 1 + 2 = 3$，$R(3) = 1 + 2(1 + 3) = 9$。 不同的子問題只有 $n(n+1)/2$ 個，指數次的呼叫
全部是重算。

\item (a) 表如下（列是 $X$，欄是 $Y$）：
{\footnotesize
\begin{center}
\begin{tabular}{cc*{8}{c}}
& &  & \texttt{C} & \texttt{H} & \texttt{E} & \texttt{A} & \texttt{T} & \texttt{E} & \texttt{R} \\
\midrule
0 &  & 0 & 0 & 0 & 0 & 0 & 0 & 0 & 0 \\
1 & \texttt{T} & 0 & $\uparrow$0 & $\uparrow$0 & $\uparrow$0 & $\uparrow$0 & $\nwarrow$1 & $\leftarrow$1 & $\leftarrow$1 \\
2 & \texttt{E} & 0 & $\uparrow$0 & $\uparrow$0 & $\nwarrow$1 & $\leftarrow$1 & $\uparrow$1 & $\nwarrow$2 & $\leftarrow$2 \\
3 & \texttt{A} & 0 & $\uparrow$0 & $\uparrow$0 & $\uparrow$1 & $\nwarrow$2 & $\leftarrow$2 & $\uparrow$2 & $\uparrow$2 \\
4 & \texttt{C} & 0 & $\nwarrow$1 & $\leftarrow$1 & $\uparrow$1 & $\uparrow$2 & $\uparrow$2 & $\uparrow$2 & $\uparrow$2 \\
5 & \texttt{H} & 0 & $\uparrow$1 & $\nwarrow$2 & $\leftarrow$2 & $\uparrow$2 & $\uparrow$2 & $\uparrow$2 & $\uparrow$2 \\
6 & \texttt{E} & 0 & $\uparrow$1 & $\uparrow$2 & $\nwarrow$3 & $\leftarrow$3 & $\leftarrow$3 & $\nwarrow$3 & $\leftarrow$3 \\
7 & \texttt{R} & 0 & $\uparrow$1 & $\uparrow$2 & $\uparrow$3 & $\uparrow$3 & $\uparrow$3 & $\uparrow$3 & $\nwarrow$4 \\
\end{tabular}
\end{center}
}
從 $[7, 7]$：$\nwarrow$ 選 \texttt{R} 到 $[6, 6]$，$\nwarrow$ 選 \texttt{E} 到 $[5, 5]$，
$\uparrow$ 到 $[4, 5]$，$\uparrow$ 到 $[3, 5]$，$\leftarrow$ 到 $[3, 4]$，$\nwarrow$ 選
\texttt{A} 到 $[2, 3]$，$\nwarrow$ 選 \texttt{E} 到 $[1, 2]$，$\uparrow$ 到 $[0, 2]$ 停。 LCS 是
\texttt{EAER}，長度 $4$。
(b) \texttt{CHER} 和 \texttt{EAER} 兩個。
(c) \texttt{CHER}。

\item $w$、$e$、$\mathit{root}$ 依區間長度：
{\footnotesize
\begin{center}
\begin{tabular}{c c p{0.55\linewidth} c c}
\toprule
$(i, j)$ & $w$ & 每個 $r$ 的成本 & $e$ & $\mathit{root}$ \\
\midrule
$(1,1)$ & 8 & $r=1$：$0+0+8=8$ & 8 & 1 \\
$(2,2)$ & 7 & $r=2$：$0+0+7=7$ & 7 & 2 \\
$(3,3)$ & 10 & $r=3$：$0+0+10=10$ & 10 & 3 \\
$(1,2)$ & 12 & $r=1$：$0+7+12=19$；$r=2$：$8+0+12=20$ & 19 & 1 \\
$(2,3)$ & 15 & $r=2$：$0+10+15=25$；$r=3$：$7+0+15=22$ & 22 & 3 \\
$(1,3)$ & 20 & $r=1$：$0+22+20=42$；$r=2$：$8+10+20=38$；$r=3$：$19+0+20=39$ & 38 & 2 \\
\bottomrule
\end{tabular}
\end{center}
}
空樹 $e[i, i-1] = 0$，$w[1,0] = 1$、$w[2,1] = 3$、$w[3,2] = 2$、$w[4,3] = 2$。 最佳樹的根是
$k_2$，左子節點 $k_1$、右子節點 $k_3$，四個外部節點 $d_0, d_1$ 掛在 $k_1$ 下、$d_2, d_3$ 掛在
$k_3$ 下，深度都是 $2$。 成本 $38$，平均 $38 / 20 = 1.9$ 次。 讓次數最多的 $k_3$ 當根的
樹成本是 $39$。 改用「走到外部節點也算一次」時多了 $q_0 + q_1 + q_2 + q_3 = 8$，
$e[1, 3] = 46$，最佳樹不變。

\item (a) $K$ 表（欄是容量 $0$ 到 $8$）：
{\footnotesize
\begin{center}
\begin{tabular}{l*{9}{c}}
$i=0$ & 0 & 0 & 0 & 0 & 0 & 0 & 0 & 0 & 0 \\
$i=1$ & 0 & 0 & 5 & 5 & 5 & 5 & 5 & 5 & 5 \\
$i=2$ & 0 & 0 & 5 & 5 & 5 & 9 & 9 & 14 & 14 \\
$i=3$ & 0 & 0 & 5 & 5 & 10 & 10 & 15 & 15 & 15 \\
$i=4$ & 0 & 0 & 5 & 10 & 10 & 15 & 15 & 20 & 20 \\
\end{tabular}
\end{center}
}
還原：$K[4, 8] = 20 \ne K[3, 8] = 15$，拿第 $4$ 件，剩 $5$；$K[3, 5] = 10 \ne K[2, 5] = 9$，拿
第 $3$ 件，剩 $1$；$K[2, 1] = K[1, 1] = K[0, 1] = 0$，第 $2$、$1$ 件不拿。 拿第 $3$、$4$ 件，重
$7$、價值 $20$。
(b) $d = \langle 0, 0, 5, 10, 10, 15, 20, 20, 25 \rangle$，$d[8] = 25$，例如第 $4$ 件拿兩次再
加第 $1$ 件。 這是無限背包的答案，不是 0/1 背包的。

\item (a) $21 \times (10^9 + 1)$，大約 $2 \times 10^{10}$ 格，就算每格只用 $4$ 個位元組也要
$80$ GB，開不出來。
(b) 從 $(20, W)$ 往下，每一層決定一件物品拿或不拿，所以走到第 $i$ 層時，$c$ 一定是 $W$ 減掉
第 $i+1..20$ 件中某些物品的重量和，最多 $2^{20 - i}$ 種。 全部加起來最多
$\sum_{i=0}^{20} 2^{20-i} = 2^{21} - 1 \approx 2.1 \times 10^6$ 個子問題。 $c$ 的值很大又很分
散，不能拿來當陣列索引，表要用第~12~章的雜湊表，以 $(i, c)$ 為鍵。

\item 表如下，$D[5, 6] = 4$：
{\footnotesize
\begin{center}
\begin{tabular}{cc*{7}{c}}
& &  & \texttt{S} & \texttt{T} & \texttt{R} & \texttt{O} & \texttt{N} & \texttt{G} \\
\midrule
0 &  & 0 & 1 & 2 & 3 & 4 & 5 & 6 \\
1 & \texttt{S} & 1 & 0 & 1 & 2 & 3 & 4 & 5 \\
2 & \texttt{P} & 2 & 1 & 1 & 2 & 3 & 4 & 5 \\
3 & \texttt{O} & 3 & 2 & 2 & 2 & 2 & 3 & 4 \\
4 & \texttt{R} & 4 & 3 & 3 & 2 & 3 & 3 & 4 \\
5 & \texttt{T} & 5 & 4 & 3 & 3 & 3 & 4 & 4 \\
\end{tabular}
\end{center}
}
依「左上、上、左」的順序還原：$[5,6]$ 是 \texttt{T} 換成 \texttt{G}，$[4,5]$ 是 \texttt{R} 換成
\texttt{N}，$[3,4]$ 是 \texttt{O} 相同，$[2,3]$ 是 \texttt{P} 換成 \texttt{R}，$[1,2]$ 是插入
\texttt{T}，$[1,1]$ 是 \texttt{S} 相同。 對齊是
\begin{center}
\texttt{S - P O R T}\\
\texttt{S T R O N G}
\end{center}
一次插入、三次替換，共 $4$ 步。

\item 把 $X_{i-1}$ 改成 $Y_j$ 的任何一種改法，最後再刪掉結尾的 $y_j$，就把 $X_{i-1}$ 改成了
$Y_{j-1}$，所以 $D[i-1, j-1] \le D[i-1, j] + 1$。 同理，先在 $X_{i-1}$ 後面插入 $x_i$ 得到
$X_i$，再用最好的方法把 $X_i$ 改成 $Y_{j-1}$，所以 $D[i-1, j-1] \le 1 + D[i, j-1]$。 $x_i = y_j$
時 $\delta = 0$，遞迴式的第一項 $D[i-1, j-1]$ 不大於另外兩項，最小值就是它。 所以這時可以直接
取左上角，和 LCS 一樣。

\item (a) $L = \langle 1, 2, 2, 3, 3, 4, 4, 5, 3 \rangle$，$\mathit{prev} = \langle \text{--}, 1,
1, 2, 3, 5, 5, 6, 3 \rangle$。 最大的是 $L[8] = 5$，走回去 $8 \to 6 \to 5 \to 3 \to 1$，LIS 是
$\langle 3, 4, 6, 10, 11 \rangle$。
(b) 每一步之後的 tail 值（括號裡是 $\mathit{prev}[i]$）：
$i = 1$：$3$（--）；$i = 2$：$3, 8$（$1$）；$i = 3$：$3, 4$（$1$）；$i = 4$：$3, 4, 12$（$3$）；
$i = 5$：$3, 4, 6$（$3$）；$i = 6$：$3, 4, 6, 10$（$5$）；$i = 7$：$3, 4, 6, 7$（$5$）；
$i = 8$：$3, 4, 6, 7, 11$（$7$）；$i = 9$：$3, 4, 5, 7, 11$（$3$）。
$\mathit{len} = 5$，$\mathit{tail}[5] = 8$，走回去 $8 \to 7 \to 5 \to 3 \to 1$，LIS 是
$\langle 3, 4, 6, 7, 11 \rangle$，和 (a) 不同，一樣長。
(c) 最後的 tail 值是 $3, 4, 5, 7, 11$，其中 $5$ 是 $A[9]$，排在 $7$、$11$ 後面，不是子序列。

\item 令 $G[i, t]$ 為前 $i$ 科共分 $t$ 小時的最高總分。 $G[0, t] = 0$；
$G[i, t] = \max_{0 \le h \le t} \bigl(G[i-1, t-h] + f_i(h)\bigr)$，第 $i$ 科分 $h$ 小時。 答案
$G[3, 4]$。
\begin{center}
\begin{tabular}{c*{5}{c}}
\toprule
$t$ & 0 & 1 & 2 & 3 & 4 \\
\midrule
$G[1, t]$ & 0 & 5 & 8 & 10 & 11 \\
$G[2, t]$ & 0 & 5 & 8 & 12 & 15 \\
$G[3, t]$ & 0 & 5 & 9 & 12 & 16 \\
\bottomrule
\end{tabular}
\end{center}
例如 $G[2, 3] = \max(10 + 0, 8 + 3, 5 + 7, 0 + 10) = 12$（第 $2$ 科分 $2$ 小時）。 $G[3, 4] =
\max(15 + 0, 12 + 4, 8 + 6, 5 + 7, 0 + 8) = 16$，第 $3$ 科分 $1$ 小時；回到 $G[2, 3]$，第 $2$
科分 $2$ 小時；剩 $G[1, 1]$，第 $1$ 科分 $1$ 小時。 最高 $16$ 分（$5 + 7 + 4$）。 一般情況有
$n(T+1)$ 格，每格最多試 $T + 1$ 個 $h$，時間 $\Theta(nT^2)$。 這也是擬多項式（$T$ 是數值）。

\item 四個矩陣如下，粗體是這一層變小的格子。
{\small
\[
D^{(0)} = \begin{pmatrix} 0 & 5 & 4 \\ 2 & 0 & -3 \\ 4 & \infty & 0 \end{pmatrix},\;
D^{(1)} = \begin{pmatrix} 0 & 5 & 4 \\ 2 & 0 & -3 \\ 4 & \mathbf{9} & 0 \end{pmatrix},\;
D^{(2)} = \begin{pmatrix} 0 & 5 & \mathbf{2} \\ 2 & 0 & -3 \\ 4 & 9 & 0 \end{pmatrix},\;
D^{(3)} = \begin{pmatrix} 0 & 5 & 2 \\ \mathbf{1} & 0 & -3 \\ 4 & 9 & 0 \end{pmatrix}.
\]
}
$d_{32}^{(1)} = 4 + 5$（$3 \to 1 \to 2$），$d_{13}^{(2)} = 5 + (-3)$（$1 \to 2 \to 3$），
$d_{21}^{(3)} = -3 + 4$（$2 \to 3 \to 1$）。

\item (a)、(c) 有，(b) 沒有。 (a)：有向無環圖上的路徑一定是簡單的。 若最長路徑經過 $v$，把
$s \leadsto v$ 那段換成更長的一段，接上 $v \leadsto t$ 那段後如果重複經過某個點 $u$，就表示
$u$ 能走到 $v$、$v$ 又能走到 $u$，形成環，矛盾；所以剪下貼上成立，可以依拓樸順序做 DP（第~22~
章）。 (b)：第~\ref{sec:dp-substructure}~節的三角形就是反例。 (c)：恰好 $k$ 條邊的最佳走法，
去掉最後一條邊 $(u, t)$ 後，一定是恰好 $k - 1$ 條邊到 $u$ 的最佳走法；點可以重複，兩段不會互
相限制。 這就是第~\ref{sec:dp-apsp}~節「限制邊數」的子問題。
\end{enumerate}
